\documentclass[11pt,a4paper]{article}
\usepackage[margin=2.2cm]{geometry}
\usepackage{amsmath,amssymb}
\usepackage{booktabs}
\usepackage{graphicx}
\usepackage{xcolor}
\usepackage{enumitem}
\usepackage{float}
\usepackage{longtable}
\usepackage{array}
\usepackage[hidelinks]{hyperref}
\usepackage{listings}
\lstset{basicstyle=\ttfamily\small,breaklines=true,columns=fullflexible}
\newcommand{\U}{\ensuremath{U}}
\newcommand{\src}[1]{{\footnotesize\texttt{#1}}}
\setlist{itemsep=2pt,topsep=3pt}
\sloppy

\title{Learning Bengali Handwriting Layout from Real Pages, v2\\[2pt]
\large Measurement-noise correction, fit-shrinkage correction, line clearance\\
and a corrected evaluation protocol for the BengaliScribe layout module}
\author{BengaliScribe / \texttt{scribe\_layout\_v2} working notes}
\date{12 October 2026}

\begin{document}
\maketitle

\begin{abstract}
BengaliScribe generates handwritten Bengali one word at a time and places the
words on a page with hand-picked random constants. Version~1 of the layout
module (\texttt{scribe\_layout/}) replaced those constants with a
page-autoregressive chain of mixture density networks (MDNs). The chain is
trained on line geometry extracted from 963 real page scans (\texttt{hcr/996}).
This report documents version~2 (\texttt{scribe\_layout\_v2/}) and is
self-contained. Version~2 changes the word anchor, which is now an inverse-noise
weighted mix of two height estimators. It estimates and removes anchor
measurement noise from the drawn wobble, and keeps neighbouring lines at least
as far apart as real lines are. It undoes the trajectory fit's shrinkage of the
wobble targets, first with a line-level table and then, after an independent
review found that table wrong for single words, with a per-word leverage
correction. It models line curvature as noise-free by extreme deconvolution,
because fitted curvature on short lines is mostly measurement noise. It also
adds a learned word rotation, places words by ink extent, resets the planner
at each page, keeps every line's ink inside the page margins, and corrects a
biased end-to-end classifier two-sample test (C2ST). The default model is
\texttt{runs/default\_unfitw\_dc}. A per-word correction calibrated on the
actual Huber trajectory fit, rather than on least squares, was tried and is
worse on all three plan-level C2STs on both test and val, so it was not
adopted.

On the held-out test split (unseen texts; 5 seeds), the leverage model reaches
a plan-level line C2ST AUC of $0.521\pm0.019$ and the default model, which adds
the curvature deconvolution, $0.535\pm0.012$ (v1: $0.566\pm0.021$). The
comparable within-page real-vs-real null is $0.496\pm0.019$, so the results
are close to the null but not at it. The design choices were compared on the
test split; a check on the validation split, run after the fact, gives
$0.527\pm0.030$ and $0.536\pm0.027$ against v1's $0.560\pm0.024$. End-to-end
(BengaliScribe render $\rightarrow$ real detector $\rightarrow$ the same
extraction), page-level C2ST stays high (0.84--0.89), and the extraction still
splits 11--13\% of the drawn lines of the default model. We report both results
and do not claim that the generated layouts are indistinguishable from real
ones.
\end{abstract}

\tableofcontents
\newpage

%=====================================================================
\section{Summary}

\paragraph{What v2 changes.} Twelve changes, each described in
Section~\ref{sec:changes} as motivation $\rightarrow$ method $\rightarrow$
parameters $\rightarrow$ evidence:
\begin{enumerate}
\item a combined word anchor with noise deconvolution (\S\ref{sec:anchor});
\item the drawn wobble de-noised by a factor $\approx0.83$ (\S\ref{sec:denoise});
\item rejection sampling for line clearance (\S\ref{sec:clear});
\item a line-level ``unfit'' table that undoes fit shrinkage (\S\ref{sec:unfit}), which had a per-word bug;
\item a per-word leverage correction that fixes that bug (\S\ref{sec:lev});
\item a curvature head that models noise-free curvature by extreme
  deconvolution, built on top of the leverage model. This is the default
  model, \texttt{runs/default\_unfitw\_dc} (\S\ref{sec:dc});
\item word rotation $\beta=0.524\times$ the line tangent angle (\S\ref{sec:tilt});
\item word placement by ink extent (\S\ref{sec:ink});
\item a planner reset at each page (\S\ref{sec:reset});
\item a margin guard that keeps each line's ink inside the top and bottom
  margins and replans a line that would cross the bottom one on a new page
  (\S\ref{sec:margin});
\item a balanced grouped C2ST for the end-to-end check (\S\ref{sec:balanced});
\item a latent writer-style model, whose context bug is only partly fixed and which is not the default (\S\ref{sec:latent}).
\end{enumerate}
The default model is \texttt{runs/default\_unfitw\_dc} (item 6). One further
change was tried and rejected: a per-word correction calibrated on the Huber
trajectory fit instead of least squares (\texttt{runs/default\_unfitwh\_dc})
is worse than the default on line, page and FFT C2ST on both test and val
(\S\ref{sec:huber}).

\paragraph{Measured results, in brief.}
\begin{itemize}
\item \emph{Plan-level} (generator output re-measured by the extraction code; test split; 5 seeds):
  line C2ST v1 $0.566\pm0.021$, v2 with the old table $0.539\pm0.016$, v2 with
  the leverage correction $0.521\pm0.019$, and the default model
  (leverage plus curvature deconvolution) $0.535\pm0.012$. Within-page null
  $0.496\pm0.019$. Page C2ST: $0.654$, $0.660$, $0.581$ and $0.590$
  respectively (Table~\ref{tab:plan}).
\item \emph{Validation split}, checked after the fact with the same protocol
  (\S\ref{sec:val}): line C2ST v1 $0.560\pm0.024$, old table $0.566\pm0.028$,
  leverage $0.527\pm0.030$, default $0.536\pm0.027$; page C2ST $0.631$,
  $0.660$, $0.575$ and $0.593$. The leverage-family models beat v1 on val too.
  The old table's test advantage over v1 on line C2ST did not replicate.
\item \emph{Controlled 2-D FFT} page-spectrum C2ST (test): v1 $0.567\pm0.034$,
  old table $0.653\pm0.024$, leverage $0.566\pm0.025$, default
  $0.550\pm0.020$. The old table's regression on this metric is what exposed
  the per-word overshoot (\S\ref{sec:fftdiag}).
\item \emph{Curvature} (\S\ref{sec:dc}): the share of real test lines with
  $n\ge5$ whose sag exceeds $0.5\,\U$ is 0.108; the leverage model gives
  0.175 and the default model 0.093. The leverage model and the default model
  differ within seed noise on every C2ST, on test and on val. The default was
  chosen for its curvature and for the halved end-to-end line splitting, not
  for C2ST.
\item \emph{End-to-end}, balanced C2ST, model with the old table
  (\texttt{e2e\_pagereset}): line C2ST v2 0.684, v2 with writer context 0.619,
  manual 0.905; page C2ST 0.834, 0.889 and 0.999. Inter-word gap $W_1$: 0.404,
  0.267 and 0.685\,\U.
\item \emph{End-to-end with the leverage model} (\texttt{e2e\_unfitw}, same
  protocol): line C2ST v2 0.699 and v2\_ctx 0.613; page C2ST 0.851 and 0.842;
  gap $W_1$ 0.376 and 0.197\,\U; pitch $W_1$ 0.346 and 0.480\,\U, about half
  the old-table values (Table~\ref{tab:e2e-balanced}).
\item \emph{End-to-end with the default model} (\texttt{e2e\_unfitw\_dc}, same
  protocol): line C2ST v2 0.719 and v2\_ctx 0.654; page C2ST 0.844 and 0.847;
  gap $W_1$ 0.387 and 0.183\,\U. The extraction splits 52 of 469 drawn v2
  lines (11\%), against 89 of 462 (19\%) with the leverage model. Slope and
  drift $W_1$ rise, but the split lines are the steep and bowed ones, so these
  $W_1$ values are not comparable across runs with different split rates
  (\S\ref{sec:e2e}).
\item \emph{Final end-to-end result} (\texttt{e2e\_final}: the default model
  plus the margin guard, same protocol; Table~\ref{tab:e2e-final}): line C2ST
  v2 0.719 and v2\_ctx 0.619; page C2ST 0.850 and 0.845; gap $W_1$ 0.425 and
  0.198\,\U. The extraction splits 57 of 453 drawn v2 lines (12.6\%). This is
  the same as \texttt{e2e\_unfitw\_dc} within noise.
\item \emph{Negative result} (\S\ref{sec:huber}): the Huber-aware per-word
  correction scores line/page/FFT C2ST $0.544/0.614/0.583$ on test and
  $0.551/0.619/0.551$ on val, against $0.535/0.590/0.550$ and
  $0.536/0.593/0.535$ for the default. It was not adopted.
\end{itemize}

\paragraph{How to read the numbers.} Lower $W_1$ and lower AUC are better. An
AUC of 0.5 means the classifier cannot separate the two samples \emph{with
these features and this amount of data}. It does not mean the samples are
indistinguishable in general. All numbers come from the run files named in
each table. Where a number comes from a comparison script whose output was not
written to a run file, the table says so.

%=====================================================================
\section{Recap of v1}
\label{sec:v1}

\subsection{Data}
\texttt{hcr/996/Images/} holds 995 RGB scans of unruled A4 pages at 300\,dpi.
Each page is a handwritten copy of one of 27 printed source texts. There are no
word boxes, no line boxes and no writer IDs. After excluding unusable pages
(remarks such as upside-down, ruled or printed; ``N/A'' text; too few words) and
a run of 28 typeset pages missed by the metadata (\texttt{image\_744}--\texttt{771}),
963 pages remain usable. The split is \emph{by source text}, with seed 0:
\begin{center}
\begin{tabular}{lrrrr}
\toprule
split & texts & pages & prose lines (v1 extraction) & inter-word gaps\\
\midrule
train & 18 & 614 & 5312 & 25\,644\\
val   & 4 \{1, 8, 20, 22\} & 144 & 1078 & 5\,000\\
test  & 5 \{4, 6, 11, 15, 18\} & 193 & 1488 & 7\,032\\
\bottomrule
\end{tabular}
\end{center}
v2 re-extracts with the combined anchor (\S\ref{sec:anchor}). Its test split has
1495 lines on the same 193 pages (\src{runs/e2e\_pagereset/results.md}).

\subsection{Geometry extraction and the writer unit \U}
The \texttt{unified\_cascade} word detector (Swin-T Mask R-CNN; configuration
\texttt{p1\_best}) supplies 50-point word polygons. Lines are rebuilt from the
words, because the detector's own line polygons merge sloped lines. Ink is
binarised with Otsu's threshold, capped at grey level 170. Each word yields its
ink extents, its headline (\emph{matra}) row with a confidence
$m_\text{frac}$, its median ink row \textit{mid}, and its core height (80th
minus 20th percentile of ink rows). Missed ink is recovered as ``orphans''.
Words are then chained greedily into lines, chain fragments are merged, words
the detector merged are split at blank-column runs, and the trailing digit
legend is tagged and excluded.

All geometry is expressed in a per-page \textbf{writer unit} $\U$: the median
word core height. On 300\,dpi scans it ranges over 31--54\,px (5th--95th
percentile), about 3.5\,mm. Among the three candidate units, $\U$ gave the most
stable ratio of line pitch to unit across writers (CV 0.189). BengaliScribe's
``body'' measures $1.114\,\U$.

\subsection{Trajectory fit}
For a line of $n$ words with centres $x_k$ and anchors $y_k$ (in \U), with
$x_m=\tfrac12(x_1+x_n)$,
\begin{equation}
y_k = c_0 + c_1\,(x_k-x_m) + c_2\,(x_k-x_m)^2 + r_k ,
\label{eq:traj}
\end{equation}
fitted by Huber IRLS (5 iterations, $c=1.345$). $c_1$ is the slope, $c_2$ the
curvature, $r_k$ the per-word residuals, and $s=1.4826\,\mathrm{MAD}(r)$ the
line's residual scale (``wobble''). The model order depends on $n$: quadratic
for $n\ge5$, linear for $n=3,4$, offset only for $n=2$. Centring at $x_m$
removes the basis collinearity between $c_1$ and $c_2$ (Spearman $-0.60\to-0.04$).
In v1 a word's anchor was its headline row when $m_\text{frac}\ge0.25$, and
otherwise \textit{mid} minus the page's median (\textit{mid}$-$headline)
offset. The same function, \texttt{dataset.measure\_page}, measures both real
and generated layouts.

On the training split the key statistics are: inter-word gap median
$1.14\,\U$, lag-1 Spearman correlation of gaps within a line 0.49; slope ICC
over pages 0.77, with consecutive-line slope correlation 0.87; residual scale
median $0.098\,\U$, ICC 0.08; residual lag-1 $-0.40$, mostly an artefact of the
fit.

\subsection{The MDN chain}
Every generator outputs, per line, a \texttt{LinePlan}
$(x_\text{start}, g_{1..n-1}, y_c, c_1, c_2, r_{1..n})$. Word $k$ is placed at
$x^{(0)}_k = x_\text{start}+\sum_{j<k}(w_j+g_j)$, with headline row
$y_c + c_1 d_k + c_2 d_k^2 + r_k$, where $d_k=x_k-x_m$. The model is
autoregressive over the page:
\begin{align}
p(\text{page}) = \prod_\ell\;
 & p(c_{1,\ell}\mid H_\ell)\,
   p(c_{2,\ell}\mid H_\ell,c_{1,\ell})\,
   p(\log(s_\ell{+}0.02)\mid H_\ell,c_{1,\ell},c_{2,\ell})\,
   p(x_{\text{start},\ell}\mid\cdot)\,
   p(\text{pitch}_\ell\mid\cdot) \nonumber\\
 & \times \prod_{k} p\big(g_{\ell,k}\mid w_{\ell,k},w_{\ell,k+1},
     x_{\ell,k},g_{\ell,k-1},\bar g_{\ell,<k},
     c_{1,\ell},c_{2,\ell},s_\ell,H^{\text{gap}}_\ell\big)
   \prod_{k} p\big(r_{\ell,k}\mid s_\ell,k/n_\ell\big).
\label{eq:factor}
\end{align}
$H_\ell$ summarises the page's earlier lines: count, mean and last slope, mean
curvature, mean log residual scale, mean and last pitch, mean line start,
gap statistics, and missing flags. Each factor is a 1-D Gaussian mixture
($K=4$) whose parameters come from a two-layer MLP (width 64, SiLU). The
weights are fitted by maximum likelihood with teacher forcing (Adam, learning
rate $10^{-3}$, weight decay $10^{-5}$, early stopping on validation NLL with
patience 30). Targets that a line cannot provide are masked out of the
likelihood. Writer-specific generation feeds a real page's measured lines in
as history. A streaming \texttt{PagePlanner}
(\texttt{begin\_line}/\texttt{next\_gap}/\texttt{end\_line}) drives
BengaliScribe's greedy line wrapping.

v1 shipped without test results: the result tables in its report were never filled in. The v1 numbers
in this report come from the v2 evaluation harness, which re-evaluates the v1
checkpoint \texttt{scribe\_layout/runs/default/model.pt}.

%=====================================================================
\section{Related work and positioning}
\label{sec:related}
We surveyed ten papers that apply Fourier methods to handwriting or to
sequence modelling (\texttt{docs/handwriting\_layout\_research\_gap.md}).
\begin{itemize}
\item \textbf{Glyph-level Fourier features} for recognition: contour Fourier
  descriptors (Lai \& Suen, 1981) and 1-D Fourier features of offline digit
  images (Liu et al., 2002). They describe single glyphs and say nothing about
  placement.
\item \textbf{Learned spectral filters on online pen trajectories} for writer
  verification: SPECTRUM (Zhang et al., ACM MM 2025). It is discriminative,
  needs online $x,y$ and pressure data, and its frequency branch helped only in
  multi-scale fusion.
\item \textbf{Frequency-aware word-image synthesis}: FW-GAN (Khoa et al.,
  2025). It works on single word crops. Its authors note that FID, KID and
  recognition scores do not capture spacing regularity.
\item \textbf{Generic spectral operators and mixers}: FNO, FNet, FNetAR and
  FourierFT. They need uniform grids or long sequences. With 4--8 gaps per
  line, a gap spectrum carries little more than a lag-1 correlation.
\item \textbf{FF-ViT} (Xu et al., 2024) regresses a single global probe angle
  from blur in endomicroscopy images. It does not model where content sits in
  the image.
\end{itemize}
None of these papers learns a generative model of page layout (line
trajectories, gap sequences, line spacing) from offline scans without pen data
or layout annotations. None models the noise of its own geometry measurements
before training a generator. None evaluates layout realism by re-measuring
generated pages with the same extractor and running distribution tests. v2
contributes to the second and third points: the noise deconvolution of
\S\ref{sec:anchor} and the re-measurement protocols of
\S\ref{sec:protocol}. Section~5 of the survey predicted where spectral methods
would be the wrong tool: short sequences, writer habit at frequency zero, and
cross-line constraints. Our 1-D and 2-D Fourier experiments
(\S\ref{sec:fourier1d}, \S\ref{sec:fft2d}) test those predictions.

%=====================================================================
\section{The v2 changes}
\label{sec:changes}

\subsection{Combined word anchor and anchor-noise deconvolution}
\label{sec:anchor}
\paragraph{Motivation.} The v1 anchor is the detected headline row. The v1
report already noted that headline misdetection puts some anchors on the wrong
row, so part of the measured wobble is measurement noise. A generator trained
on the raw residuals reproduces that detector noise as if it were handwriting.

\paragraph{Method.} Every word has two height estimators, and a real vertical
shift $t$ of the word moves both:
$A = t + e_A$ is the headline row, and $C = t + e_C$ is the median ink row
minus the page's median (\textit{mid}$-$headline) offset. If $e_A$ and $e_C$
are independent of each other and of $t$, the residuals $r_A,r_C$ of the
trajectory fit~\eqref{eq:traj}, applied to each estimator separately, satisfy
\begin{equation}
\operatorname{cov}(r_A,r_C)=\sigma_t^2,\qquad
\operatorname{var}(r_A)-\sigma_t^2=\sigma_A^2,\qquad
\operatorname{var}(r_C)-\sigma_t^2=\sigma_C^2 .
\end{equation}
The v2 anchor is the inverse-noise weighted combination
\begin{equation}
a = w\,A + (1-w)\,C,\qquad
w=\frac{\sigma_A^{-2}}{\sigma_A^{-2}+\sigma_C^{-2}},\qquad
\sigma_\text{comb}^2=\big(\sigma_A^{-2}+\sigma_C^{-2}\big)^{-1}.
\end{equation}
Words with $m_\text{frac}<0.25$ use $C$ alone
(\texttt{extract.anchor}).

\paragraph{Parameters and how they were estimated.} The estimator is
\texttt{anchor\_noise.py}. It uses train-split lines with at least 5
headline-confident words, 21\,880 words in all, and winsorises residuals at
the 1st and 99th percentiles (\src{runs/anchor\_noise.json}). It gives
$\sigma_t^2=0.00906\,\U^2$, $\sigma_A^2=0.01176$, $\sigma_C^2=0.00631$,
$w=0.349$ and $\sigma_\text{comb}^2=0.00411$. The headline-only noise share
is $\sigma_A^2/\operatorname{var}(r_A)=0.565$. For the combined anchor,
the share is $0.00411/(0.00906+0.00411)=0.31$ (computed from the same file).

\paragraph{Evidence and caveats.} The independence assumption does not hold
exactly (review M3, \S\ref{sec:review}). The reviewer added two more estimators,
body bottom and a longest-run headline, and the one-factor tetrad test fails:
errors from the same mechanism correlate. Across the plausible values of
$\sigma_t^2$, the \emph{headline-only} noise share lies anywhere between 38\% and
66\%. The winsorising level also moves it: 0.61 without clipping, 0.526 at
5\%. We therefore do not report ``57\% of the wobble is noise'' as a finding.
The \emph{combined-anchor} share, which is what the generator uses, stays at
0.29--0.33 across clip levels.

\subsection{De-noised wobble at generation time}
\label{sec:denoise}
\paragraph{Motivation.} The model learns the \emph{measured} per-word wobble,
which includes anchor noise. The real handwriting should only receive the
physical part. Generated pages are re-measured by the same extractor, which adds
its own noise back on top.

\paragraph{Method.} Drawn residuals are multiplied by
\begin{equation}
f_\text{wobble}=\sqrt{1-\frac{\sigma^2_\text{comb}}{\sigma_t^2+\sigma^2_\text{comb}}}
\end{equation}
(\texttt{bscribe\_layout.wobble\_factor}). Fitting shrinks noise and true
wobble by the same factor under the i.i.d.\ assumption, so the noise share is
the same at every line length. The review checked that this does not
double-count.

\paragraph{Parameters.} $f_\text{wobble}=\sqrt{1-0.31}\approx0.83$, derived
from the train-split values of \S\ref{sec:anchor}. Over the combined-share
range 0.29--0.33 the factor is 0.82--0.84, so the M3 uncertainty barely
affects it.

\paragraph{Evidence.} Fig.~\ref{fig:fixes} (middle panel) shows the same
generated words with the factor applied. The effect is subtle at page scale.
There is no quantitative isolation of this step alone. The FFT diagnostic of
\S\ref{sec:fftdiag} found $\times0.83$ (0.658) better than $\times1.0$
(0.683), but that was measured on the old-table model.

\subsection{Line clearance by rejection sampling}
\label{sec:clear}
\paragraph{Motivation.} The model samples each line's slope and curvature
from the page history without regard to the previous line's trajectory, so
neighbouring lines can converge or touch. Real neighbouring lines almost never
do. Without rejection, 9--15\% of generated neighbouring line pairs violate the
real 1st-percentile clearance, against 1.4\% for real pairs
(\src{runs/eval/ablation.md}, \texttt{clearance\_violation\_rate}).

\paragraph{Method.} For consecutive lines that overlap by at least $5\,\U$ in
$x$, let $d(x)$ be the vertical distance between their trajectories, sampled at
40 points over the overlap. A new line is accepted if
\begin{equation}
\min_x d(x)\ \ge\ d_\text{min}\quad\text{and}\quad
\frac{\min_x d(x)}{\overline{d(x)}}\ \ge\ \rho_\text{min}.
\end{equation}
Otherwise the line's slope, curvature and log residual scale are redrawn
(\texttt{planner.resample\_shape}), up to 30 times. If no draw passes, the line
is pushed down just far enough to pass.

\paragraph{Parameters.} $d_\text{min}=2.603\,\U$ and $\rho_\text{min}=0.797$
are the 1st percentiles over 3\,740 consecutive real line pairs (both lines
$n\ge3$) on the train split. The val split gives 2.613 and 0.788
(\src{runs/line\_clearance.json}). The cap of 30 redraws was set by hand.

\paragraph{Evidence.} In the \texttt{e2e\_unfitw} run log, over 20 generated
documents (v2 and v2\_ctx, 5 writers $\times$ 2 seeds), 179 of 941 lines
(19\%) needed at least one redraw, and 61 (6.5\%) were pushed down after 30
failed redraws. Per document the counts were 3--19 redrawn and 0--9 pushed,
out of 41--55 lines.
The plan-level evaluation in Table~\ref{tab:plan} does \emph{not} apply the
rejection, which runs only inside BengaliScribe's layout. Fig.~\ref{fig:fixes}
(right) shows the result on the page.

\subsection{Undoing the fit shrinkage, attempt 1: a line-level table}
\label{sec:unfit}
\paragraph{Motivation.} The training targets $r_k$ and $s$ are residuals of the
fit~\eqref{eq:traj}, which absorbs part of each word's offset. When a model
reproduces those residuals and the evaluation re-fits the generated anchors,
the wobble is shrunk \emph{twice}. For a 5-word line, the re-measured scale
drops to about 0.6 of the sampled one.

\paragraph{Method.} By simulation (\texttt{unfit.py}; 4000 lines of i.i.d.\
$\mathcal N(0,1)$ offsets per $n$, fixed seed, the same \texttt{fit\_line}):
\begin{equation}
\text{ratio}(n)=\frac{\operatorname{median}\,s(\text{fit residuals})}
{\operatorname{median}\,s(\text{offsets})},\qquad
\text{inflation}(n)=1/\text{ratio}(n).
\end{equation}
Training targets become $\log(s\cdot\text{inflation}(n)+0.02)$ for the scale
and $r_k\cdot\text{inflation}(n)$ for every per-word residual
(\texttt{configs}: \texttt{unfit\_targets: true}). The model then learns
pre-fit (``physical'') wobble directly.

\paragraph{Parameters.} The table depends on no data
(\src{runs/unfit\_ratio.json}):
\begin{center}
\begin{tabular}{lcccccccccc}
\toprule
$n$ & 3 & 4 & 5 & 6 & 7 & 8 & 9 & 10 & 11 & 12\\
inflation$(n)$ & 2.505 & 1.52 & 3.159 & 2.076 & 1.644 & 1.408 & 1.351 & 1.29 & 1.256 & 1.234\\
\bottomrule
\end{tabular}
\end{center}
The table is non-monotone: it jumps where the model order rises at $n=3$ and $n=5$.

\paragraph{Evidence.} Plan-level line C2ST fell from $0.638\pm0.006$
(\texttt{v2\_history}, fitted targets) to $0.539\pm0.016$
(\texttt{v2\_history\_unfit}) (\src{runs/eval/ablation.md}, 5 seeds). On the
validation split the old-table model does not beat v1 on line C2ST
(\S\ref{sec:val}).
The independent review then showed that the table is wrong for single words
(M2). It matches the median robust \emph{scale} of a line, but it was also
applied to every per-word residual target. For $n=5$ the inflated targets have
a per-word RMS of 2.07 times the physical $\sigma$, while their MAD is correct.
The C2ST gain is therefore partly an artefact: the evaluation measures the same
median-of-MAD statistic that the table was calibrated on. The coordinator
reproduced the overshoot by simulation (Table~\ref{tab:levsim}).

\subsection{Undoing the fit shrinkage, attempt 2: per-word leverage}
\label{sec:lev}
\paragraph{Method.} For ordinary least squares with design matrix $V$ (centred
powers of $x$, of degree 2 if $n\ge5$, 1 if $n\ge3$, else 0), the hat matrix is
$H=V(V^\top V)^{-1}V^\top$. An i.i.d.\ offset of standard deviation $\sigma$
then leaves residual $i$ with standard deviation $\sigma\sqrt{1-h_i}$, where
$h_i=H_{ii}$ is that word's leverage. The per-word correction is
\begin{equation}
\tilde r_i = r_i\cdot\min\!\Big(\frac{1}{\sqrt{1-h_i}},\,4\Big)
\end{equation}
(\texttt{unfit.word\_inflation}). It uses the line's real word positions, and
restores standard deviation $\sigma$ at every $n$ and every position. The cap of
4 guards words that nearly fix the fit. The line-scale target
$\log(s\cdot\text{inflation}(n)+0.02)$ keeps the table, which is correct for
that statistic. Configuration: \texttt{configs/default\_unfitw.yaml}
(\texttt{unfit\_words: true}); checkpoint
\texttt{runs/default\_unfitw/model.pt}.

\paragraph{Parameters.} The leverages $h_i$ are computed per line from that
line's word centres. No data-fitted constant is involved. The MDN weights are
retrained on the new targets: best epoch 152, stopped at 183
(\src{runs/default\_unfitw/metrics.json}).

\paragraph{Evidence.} Simulation (FACTS item 4): i.i.d.\ wobble with standard
deviation 0.150 is placed at real train-line word positions and refitted with
\texttt{dataset.fit\_line}.
\begin{table}[H]\centering\small
\begin{tabular}{lcccccccccc}
\toprule
$n$ & 3 & 4 & 5 & 6 & 7 & 8 & 9 & 10 & 11 & 12\\
\midrule
raw fit residual sd & .087 & .107 & .098 & .109 & .115 & .121 & .124 & .126 & .129 & .132\\
leverage-corrected  & .151 & .149 & .142 & .147 & .148 & .150 & .150 & .149 & .149 & .149\\
old table           & .219 & .163 & .311 & .227 & .190 & .170 & .168 & .163 & .162 & .159\\
\bottomrule
\end{tabular}
\caption{Per-word standard deviation recovered from i.i.d.\ wobble with true
sd 0.150\,\U. The leverage correction recovers 0.142--0.151 at every $n$. The
old table overshoots by up to $2.1\times$ (at $n=5$).}
\label{tab:levsim}
\end{table}
At plan level the leverage model has line C2ST $0.521\pm0.019$ and page C2ST
$0.581\pm0.026$. The residual-scale $W_1$ for short lines falls from 0.039 to
0.009 ($n=5$) and from 0.028 to 0.011 ($n=6$--$7$). The controlled 2-D FFT
C2ST returns to v1's level: $0.566\pm0.025$ against $0.653\pm0.024$ for the
old table (Table~\ref{tab:plan}).

\paragraph{What it does not fix.} The leverage correction assumes i.i.d.\
wobble. Real wobble is serially correlated, with residual lag-1 $-0.38$ on the
test split against about $-0.48$ for i.i.d.\ generation
(\src{runs/eval/ablation.md}). Real residual scale also grows with line length
faster than fit shrinkage explains (review M2). The long-line bins still differ
from real ($W_1$ 0.041 at $n=8$--$10$, 0.062 at $n\ge11$).

\subsection{De-noised curvature (extreme deconvolution)}
\label{sec:dc}
\paragraph{Motivation: why generated lines split end to end.} BengaliScribe
writes the drawn word and line polygons of every generated page, so extracted
lines can be matched to the lines that were actually drawn
(\src{runs/compare\_unfitw/split\_diag.py}; an extracted piece belongs to the
drawn line whose word polygon contains its centre). In the end-to-end run of
the leverage model (\texttt{e2e\_unfitw}, method v2, 5 writers $\times$ 2
seeds), 89 of 462 drawn lines (19\%) were split into two or more extracted
lines, and 127 extracted lines were single pieces. For the manual layout it was
6 of 355 (1.7\%). At the 113 split points
(\src{split\_why.py}), the step in median ink row between the last word of one
fragment and the first word of the next has median $1.2\,\U$ (10th percentile
0.72, 90th 1.83). The line chainer joins neighbours only within
\texttt{JOIN\_DEV} $=1.0\,\U$; 74\% of the steps exceed it, and for 86\% the
fragment-merge residual is finite but at least $0.6\,\U$. The steps are not a
glyph effect: the (\textit{mid}$-$headline) offset of generated words has an
interquartile range of $0.19\,\U$, narrower than the real $0.30\,\U$.

The steps come from lines that bow too much. Let the \emph{sag} of a line be
$c_2\,(L/2)^2$ in \U, where $L$ is the distance between its first and last
word centres. At plan level on the test split (seed 0, lines with $n\ge5$;
\src{sag.py}), the share of lines with $|\text{sag}|>0.5\,\U$ is 0.108 for
real lines, 0.170 for v1 and 0.175 for the leverage model
(Table~\ref{tab:sag}). The leverage model's sag sd is $0.358$ and
$0.509\,\U$ on lines of 5--6 and 7--8 words, against $0.265$ and $0.322\,\U$
for real lines.

\paragraph{Interpretation: curvature noise.} The curvature head is sampled
before the line's word count is known (Eq.~\ref{eq:factor}), so it learns one
distribution of \emph{fitted} $c_2$ pooled over all line lengths. On lines
with 5--6 words, the fitted $c_2$ is mostly the trajectory fit's response to
per-word wobble and anchor noise. The noise model below puts numbers on this:
on the train split, the measured curvature sd falls from 0.00129 ($n=5$--6) to
0.00084--0.00092 for longer lines, while the noise-corrected sd stays at
0.0007--0.0009 for every $n$ (Table~\ref{tab:curvnoise}). True curvature looks
like a habit of similar size at every line length. The noise of short lines,
drawn on lines of any length, then becomes real bow.

\paragraph{Method.} Write the fitted curvature of line $i$ as the true value
plus fit noise,
\begin{equation}
\hat c_{2,i} = c_{2,i} + \varepsilon_i,\qquad
\varepsilon_i\sim\mathcal N(0,v_i),\qquad
v_i = \bar s^2_{n_i}\,\big[(V_i^\top V_i)^{-1}\big]_{22},
\label{eq:curvnoise}
\end{equation}
where $V_i$ is the line's centred quadratic design matrix (columns $1$,
$x-x_m$, $(x-x_m)^2$) and $\bar s^2_n$ is the per-word residual variance for
lines with $n$ words. This is the ordinary-least-squares variance of $c_2$; it
ignores the Huber weights of the actual fit and treats the words as
independent. A single line's own estimate $\text{RSS}/(n-3)$ has only $n-3$
degrees of freedom, so $\bar s^2_n$ is the mean of $\text{RSS}/(n-3)$ over all
train lines with the same $n$, with $n\ge11$ pooled into one group
(\texttt{features.curv\_noise\_var}). Following extreme deconvolution (Bovy,
Hogg and Roweis, 2011~\cite{bovy2011}), the MDN curvature head with mixture
parameters $(\pi_k,\mu_k,\sigma_k)$ then describes the \emph{noise-free}
curvature, and each measured target is scored under the mixture convolved with
its own noise:
\begin{equation}
p(\hat c_{2,i}\mid H_i,c_{1,i}) =
\sum_{k=1}^{K}\pi_k\,\mathcal N\!\big(\hat c_{2,i};\;\mu_k,\;\sigma_k^2+v_i\big)
\label{eq:xd}
\end{equation}
(\texttt{nn\_utils.MDN.nll(noise\_var)}, \texttt{model.head\_nll}). At
generation time the curvature is drawn from
$\sum_k\pi_k\,\mathcal N(\mu_k,\sigma_k^2)$, without the noise term. The
re-measurement of a generated line adds its own fit noise back, as it does for
real lines. Configuration: \texttt{configs/default\_unfitw\_dc.yaml}, which is
\texttt{default\_unfitw} plus \texttt{deconv: [curv]}; checkpoint
\texttt{runs/default\_unfitw\_dc/model.pt}. This is the default model of
\texttt{bscribe\_run.py} and \texttt{relayout.py}.

\paragraph{Learned and estimated quantities.} The per-$n$ pooled per-word
variances $\bar s^2_n$ are estimated on the train split; the design factor
$[(V_i^\top V_i)^{-1}]_{22}$ is computed per line from its word centres. Over
the 4050 train curvature targets, the mean noise variance $v_i$ is 38.4\% of
the variance of the measured curvature target. The MDN weights are retrained
with the same seed and schedule. All heads other than curvature come out
bit-identical to \texttt{default\_unfitw}; only the curvature head's
likelihood changed.

\begin{table}[H]\centering\small
\begin{tabular}{lcccc}
\toprule
word count $n$ & 5--6 & 7--8 & 9--10 & $\ge11$\\
\midrule
measured curvature sd       & 0.00129 & 0.00092 & 0.00084 & 0.00087\\
mean noise sd $\sqrt{\bar v}$ & 0.00092 & 0.00053 & 0.00053 & 0.00047\\
implied true sd             & 0.00091 & 0.00075 & 0.00066 & 0.00073\\
\bottomrule
\end{tabular}
\caption{Curvature noise decomposition on the train split
(\src{runs/compare\_unfitw/curvnoise2.py}; output printed, not saved). Implied
true sd $=\sqrt{\operatorname{var}(\hat c_2)-\bar v}$. Curvature is in
$\U^{-1}$.}
\label{tab:curvnoise}
\end{table}

\paragraph{Evidence.} Table~\ref{tab:sag} re-measures generated plans on the
test split. The deconvolved model brings the share of strongly bowed lines and
the sag spread to the real values, and its measured curvature sd is within
about 0.0001 of real for $n\le8$. For $n\ge9$ it is \emph{below} real (0.00073
vs 0.00112 at $n=9$--10, 0.00071 vs 0.00086 at $n\ge11$); these bins hold
only 76 and 34 test lines. The share of lines whose largest step between
adjacent word anchors exceeds $1\,\U$ falls from 0.169 to 0.138 (real 0.099).
Over 5 seeds at plan level (\src{runs/compare\_unfitw/drift.py}; output
printed, not saved), the curvature $W_1$ to real is 0.0002 for the leverage
model and 0.0001 for the deconvolved one, while slope (0.0045 vs 0.0044),
drift (0.158 vs 0.155) and pitch $W_1$ (0.109 vs 0.114) barely move. The
validation NLL of the curvature head, scored on the measured curvature, is
$-1.004$ nats against $-0.975$ for \texttt{default\_unfitw}. Both numbers are
predictive densities of the same measured target, so they can be compared; the
deconvolved model's NLL uses the convolved likelihood~\eqref{eq:xd}. End to
end, the split rate of drawn v2 lines falls from 19\% to 11\%
(\S\ref{sec:e2e}). Plan-level C2ST does not separate the two models: line
C2ST $0.535\pm0.012$ against $0.521\pm0.019$ on test and $0.536\pm0.027$
against $0.527\pm0.030$ on val (\S\ref{sec:val}). The deconvolved model is the
default because of the curvature statistics and the end-to-end splitting, not
because of C2ST.

\begin{table}[H]\centering\small
\setlength{\tabcolsep}{4pt}
\begin{tabular}{lcccccccc}
\toprule
 & $|\text{sag}|>0.5\,\U$ & \multicolumn{2}{c}{sag sd (\U)} & \multicolumn{4}{c}{measured curvature sd} & max.\ adj.\\
 & share, $n\ge5$ & $n=5$--6 & 7--8 & $n=5$--6 & 7--8 & 9--10 & $\ge11$ & step $>1\,\U$\\
\midrule
real (test)          & 0.108 & 0.265 & 0.322 & 0.00125 & 0.00083 & 0.00112 & 0.00086 & 0.099\\
v1                   & 0.170 & ---   & ---   & ---     & ---     & ---     & ---     & 0.173\\
v2 leverage          & 0.175 & 0.358 & 0.509 & 0.00146 & 0.00113 & 0.00105 & 0.00113 & 0.169\\
\textbf{v2 deconvolved} & 0.093 & 0.260 & 0.314 & 0.00117 & 0.00077 & 0.00073 & 0.00071 & 0.138\\
\bottomrule
\end{tabular}
\caption{Line bow at plan level, test split, seed 0, generated plans
re-measured by \texttt{measure\_page} (\src{runs/compare\_unfitw/sag.py},
which takes the v2 checkpoint as its argument; output printed, not saved).
Sag $=c_2(L/2)^2$. v2 leverage is \texttt{default\_unfitw}; v2 deconvolved
is \texttt{default\_unfitw\_dc}, the default. The last column is the share of
lines whose largest step between adjacent word anchors exceeds $1\,\U$. The
bins $n=9$--10 and $n\ge11$ hold 76 and 34 lines. v1 is reported only for the
two whole-population columns.}
\label{tab:sag}
\end{table}

\paragraph{What it does not fix.} The noise model is approximate: it ignores
the Huber weights, assumes independent words although real wobble is
serially correlated (\S\ref{sec:lev}), and pools $\bar s^2_n$ over all lines
of a length, so a line with unusually large wobble gets too small a $v_i$. The
under-dispersion for $n\ge9$ is consistent with the curvature head still not
seeing the line length, but this was not tested, and the bins are small.
Splitting is reduced, not removed (\S\ref{sec:e2e}).

\subsection{Huber-aware per-word correction (negative result, not adopted)}
\label{sec:huber}
\paragraph{Motivation.} The leverage factor $1/\sqrt{1-h_i}$ of
\S\ref{sec:lev} is exact for an ordinary-least-squares fit. The trajectory
fit \texttt{dataset.fit\_line} is a Huber IRLS fit, whose re-weighting shrinks
the residuals of high-leverage words further. A new unit test
(\src{tests/test\_unfit.py}) shows the consequence: with i.i.d.\ Gaussian
offsets, the leverage-corrected residuals of the first and last words of
lines with $n=4$--10 words stay about 10--25\% below the drawn standard
deviation. The leverage model therefore under-corrects the edge words.

\paragraph{Method.} \texttt{unfit.word\_inflation\_huber} measures the factor
on the actual fit. For each word count $n$ and word position, the factor is
$\operatorname{sd}(\text{offset})/\operatorname{sd}(\text{Huber-fit
residual})$, estimated by passing 4000 i.i.d.\ $\mathcal N(0,1)$ offset
vectors through \texttt{dataset.fit\_line} with evenly spaced words. The cap
of 4 is kept, lines longer than 30 words interpolate the $n=30$ row over
relative position, and the table is cached in
\src{runs/unfit\_word\_huber.json}. Unlike the leverage factor, it depends
on word position only, not on the line's actual word spacing. The model
\texttt{runs/default\_unfitwh\_dc} (\texttt{configs/default\_unfitwh\_dc.yaml}:
\texttt{unfit\_words: huber}, \texttt{deconv: [curv]}) differs from the
default only in this per-word factor; the line-scale target keeps the table.

\begin{table}[H]\centering\small
\begin{tabular}{llll}
\toprule
$n$ & & Huber-aware (edge \ldots\ edge) & leverage (edge \ldots\ edge)\\
\midrule
5  & all words & 4.0 (cap), 1.14, 1.46, 1.14, 4.0 (cap) & 2.96, 1.26, 1.39, 1.26, 2.96\\
6  & all words & 2.91, 1.13, 1.20, 1.21, 1.13, 2.89 & 2.37, 1.20, 1.26, 1.26, 1.20, 2.37\\
8  & edges / interior & 2.11, 2.13 / 1.10--1.17 & 1.85, 1.85 / 1.14--1.18\\
10 & edges / interior & 1.76, 1.74 / 1.08--1.15 & 1.62, 1.62 / 1.11--1.18\\
12 & edges / interior & 1.55, 1.57 / 1.06--1.16 & 1.49, 1.49 / 1.09--1.18\\
\bottomrule
\end{tabular}
\caption{Per-word inflation factors for evenly spaced words. Huber-aware:
\src{runs/unfit\_word\_huber.json}. Leverage: \texttt{unfit.word\_inflation}
evaluated on the same evenly spaced positions for this table. The Huber-aware
factors are larger at the line ends and slightly smaller inside the line. The
small left--right asymmetry of the Huber-aware values is simulation noise.}
\label{tab:huberfac}
\end{table}

\paragraph{Evidence.} The plan-level protocol of Table~\ref{tab:plan}, 5 seeds,
on test and on val (\src{runs/compare\_unfitw/compare\_5seed\_wh\_test.py}
and \src{compare\_5seed\_wh\_val.py}, outputs in the matching \texttt{.txt}
files):
\begin{table}[H]\centering\small
\begin{tabular}{llcccc}
\toprule
 & & & & FFT C2ST & resid.\ scale $W_1$\\
split & per-word factor & line C2ST & page C2ST & (contr.) & $n=5$/6--7/8--10/11+\\
\midrule
test & \textbf{leverage (default)} & $0.535\pm0.012$ & $0.590\pm0.024$ & $0.550\pm0.020$ & .007 / .011 / .042 / .061\\
     & Huber-aware & $0.544\pm0.012$ & $0.614\pm0.011$ & $0.583\pm0.029$ & .009 / .014 / .043 / .068\\
\midrule
val  & \textbf{leverage (default)} & $0.536\pm0.027$ & $0.593\pm0.031$ & $0.535\pm0.044$ & .011 / .011 / .035 / .069\\
     & Huber-aware & $0.551\pm0.022$ & $0.619\pm0.043$ & $0.551\pm0.049$ & .010 / .013 / .041 / .073\\
\bottomrule
\end{tabular}
\caption{Huber-aware against leverage per-word correction, both with the
curvature deconvolution: \texttt{runs/default\_unfitwh\_dc} against the
default \texttt{runs/default\_unfitw\_dc}. 5 seeds, mean $\pm$ sd
(\src{runs/compare\_unfitw/compare\_5seed\_wh\_test.txt} and \src{compare\_5seed\_wh\_val.txt}).}
\label{tab:huber}
\end{table}

\emph{Measured:} the Huber-aware model is worse on all three C2STs on both
splits, six of six comparisons in the same direction. Each single difference
is small, from $+0.009$ (test line C2ST) to $+0.033$ (test FFT C2ST), or
about 0.3 to 2 seed-sd depending on the comparison and on which model's sd is
used; the evidence is the consistent direction, not any one gap. The
residual-scale $W_1$ is higher in 7 of 8 bins (lower only at $n=5$ on val). Both splits had already been used (\S\ref{sec:val}), so this is a
development comparison; since the candidate was rejected, it does not
flatter the default's numbers.

\emph{Interpretation (not verified further):} the Huber-aware factor is exact
only for i.i.d.\ wobble. Real wobble is smoother than i.i.d.: the residual
lag-1 correlation is $-0.38$ on real lines against about $-0.48$ for i.i.d.\
generation (review M2, \S\ref{sec:lev}). The edge residuals of real lines also
carry line-end detection noise. Restoring the full i.i.d.\ edge variance,
with factors up to the cap of 4 at $n=5$, then probably overshoots at the
line ends, and the leverage factor's smaller edge inflation happens to sit
closer to what real lines need. Neither part of this explanation was tested
on its own.

\paragraph{Decision and open fix.} The default stays
\texttt{runs/default\_unfitw\_dc}. The leverage factor is known to be
inexact; it is kept because it scores better, not because it is right. The
proper fix, not attempted, is to calibrate the per-word correction under a
correlated-wobble model fitted to the real lag correlations, with the same
Huber fit, rather than under i.i.d.\ offsets.

\subsection{Word rotation along the trajectory}
\label{sec:tilt}
\paragraph{Motivation.} BengaliScribe pastes words unrotated, so on a curved or
sloped line each word's headline stays horizontal while the line moves. v1
measured word tilt but did not model it.

\paragraph{Method.} Each word is rotated about its ink centre at the headline
row (\texttt{write\_page.rotate\_word}), by
\begin{equation}
\theta_k=\beta\cdot\arctan\!\big(c_1+2c_2 d_k\big).
\end{equation}

\paragraph{Parameters.} $\beta=0.524$, a Huber IRLS regression of measured
headline tilt on the line tangent angle over 26\,653 train words. On val,
$\beta=0.531$ over 5\,237 words (\src{runs/default\_unfitw/tilt.json}). The
fitted intercept ($-0.009$\,rad) is not applied.

\paragraph{Evidence.} The review checked the pivot and the sign convention. No
metric isolates rotation: the line features do not include word tilt.

\subsection{Ink-extent placement}
\label{sec:ink}
\paragraph{Motivation.} The model's gaps are measured from ink edge to ink edge.
BengaliScribe word images carry blank margins, so placing words by canvas width
made every gap about $0.5\,\U$ too wide.

\paragraph{Method.} Gaps are applied between the ink extents of each word image,
using columns darker than \texttt{INK\_LEVEL} $=0.35$
(\texttt{bscribe\_layout.ink\_extents}).

\paragraph{Parameters.} \texttt{INK\_LEVEL} was set by hand to match the real
binarisation, where the Otsu cap of 170 keeps grey values darker than about
0.33. On generated pages Otsu gives about 0.39, so right ink edges come out
$0.063\,\U$ short and every measured gap is about $+0.067\,\U$ too wide
(review minor 2). This is not fixed.

\paragraph{Evidence.} End-to-end, with the biased C2ST and the same model
(\texttt{e2e\_unfit} $\rightarrow$ \texttt{e2e\_unfit\_ink}): the v2 gap $W_1$
fell from 0.602 to 0.471, and v2 with writer context from 0.381 to 0.207
(Table~\ref{tab:e2e-biased}).

\subsection{Planner reset per page}
\label{sec:reset}
\paragraph{Motivation.} Training builds the page history afresh for every page.
BengaliScribe, however, built one \texttt{PagePlanner} per document, and its
history ran across 3--5 pages. The history count reached 56--72 lines, against
at most 16 in training. History-conditioned slopes random-walked: one page
mean was $-0.158$, reaching the $-0.2$ clamp, and 2 words were drawn past the
page bottom (review C3).

\paragraph{Method.} Each new page restarts from one fixed reference state:
the real context page if one is given, else the state at the end of the
first generated page (\texttt{bscribe\_layout.layout\_v2}, \texttt{snapshot}).
All pages of a document therefore share one writer without drifting.

\paragraph{Evidence.} After the fix, no words fall past the bottom. Page slope
means stay near the writer's: for image\_104 with context they are $-0.053$,
$-0.054$, $-0.074$, $-0.009$ and $-0.016$, against $-0.023$ on the real page.

\subsection{Margin guard}
\label{sec:margin}
\paragraph{Motivation.} Found while writing the first tests for
\texttt{bscribe\_layout} (\S\ref{sec:review}). The page-break check looked
only at the line centre: a new page started when
$\text{center}+1.6\cdot\text{body}>H-m_b$. The line's slope, bow and the
words' own ink extent were ignored. A sloped or bowed first line could rise
about 1--1.4\,\U\ into the top margin, and a slope could carry ink below the
bottom margin.

\paragraph{Method.} In \texttt{bscribe\_layout.layout\_v2},
\texttt{ink\_rows} gives the top and bottom ink rows of every generated word
image (same \texttt{INK\_LEVEL} as \S\ref{sec:ink}). \texttt{ink\_span}
returns the line's ink top and bottom along its current trajectory (slope and
curvature, after the clearance step of \S\ref{sec:clear}), widened by
\texttt{RESID\_MARGIN} for the per-word wobble that is drawn afterwards.
\begin{itemize}
\item \emph{Bottom margin.} If the line's ink would pass the bottom margin
  (or the old centre test fires), the line starts a new page and is
  \emph{replanned from scratch}: its words go back to the queue, and the
  discarded draw never enters the page history. The new page therefore starts
  from the reference state of \S\ref{sec:reset}, not from a state that already
  contains the rejected line.
\item \emph{Top margin.} A page's first line is pushed down until its ink
  clears the top margin.
\end{itemize}

\paragraph{Parameters.} \texttt{RESID\_MARGIN} $=0.3\,\U$, set by hand to
about 2 sd of the drawn wobble. Wobble beyond it, and small rotation effects,
can still cross a margin.

\paragraph{Evidence.} \texttt{tests/test\_bscribe\_layout.py::TestWordsStayInsideMargins}
lays out synthetic word images over 2 seeds $\times$ (no context, context) and
asserts that every word's ink lies within the margins $\pm0.5\,\U$. It passes
on the new code and \emph{fails on the old code}, where the highest ink row
is at 178\,px against a limit of 216\,px. End to end, the guard changes no
metric beyond noise (\texttt{e2e\_final}, \S\ref{sec:e2e}).

\subsection{Balanced grouped C2ST for the end-to-end check}
\label{sec:balanced}
\paragraph{Motivation.} In the end-to-end check every page is its own group and
holds a single label. Grouped $k$-fold then changes the class balance between
folds, the logistic intercept absorbs it, and pooling the out-of-fold scores
pushes the AUC below 0.5. The old end-to-end real-vs-real ``floor'' was 0.34
(line) and 0.115 (page), which shows the protocol was broken (review C1).

\paragraph{Method.} Groups are dealt into folds separately per label, and the
reported AUC is the mean of the per-fold AUCs
(\texttt{evaluate.grouped\_kfold\_auc(balanced=True)}; end-to-end only). The
plan-level C2ST is unchanged, because there every group holds both labels.

\paragraph{Evidence.} The new real-vs-real floor over 5 seeds is 0.503 (line)
and 0.527 (page). In the \texttt{e2e\_pagereset} run it is 0.534 and 0.508
(\src{runs/e2e\_pagereset/results.md}). C2ST values from before and after this
change are \emph{not comparable}. Tables~\ref{tab:e2e-biased}
and~\ref{tab:e2e-balanced} are kept separate for this reason.

\subsection{Latent writer-style model (not the default; unresolved)}
\label{sec:latent}
\paragraph{Motivation.} The history model represents a writer only through
running means. An alternative is a writer-style vector $z$ that is fixed for the
whole page (\texttt{style: latent}). The line heads are conditioned on $z$ plus
a short previous-line term, and a small autoregressive MDN density over $z$ draws
it when no context is given.

\paragraph{Bug and partial fix.} During training, $z$ is a leave-one-out
summary over almost the whole page. With writer context, $z$ comes from half a
page. This mismatch in the precision of $z$ hurt context generation. A
causal-prefix $z$ (\texttt{runs/latent\_unfit\_fix}) improved latent+context:
slope $W_1$ $0.00977\to0.00406$, drift $W_1$ $0.262\to0.106$, line C2ST
$0.561\to0.547$. For comparison, history+context gives 0.00332, 0.110 and
0.540. The same fix \emph{regresses} the latent model without context:
slope $W_1$ $0.00533\to0.01283$ and page C2ST $0.684\to0.829$
(\src{runs/eval/ablation.md} vs \src{runs/eval\_latentfix/ablation.md}). The
likely cause is that one fixed $z$ is reused at every line position, while the
precision of $z$ in training depends on the position. This is unresolved, and
the latent model is not used.

%=====================================================================
\section{Learned and estimated parameters}
\label{sec:params}
Table~\ref{tab:params} lists every constant that shapes the v2 layout,
together with its value and where it comes from. Extraction thresholds that
define \emph{what is measured} are as in v1 and are not repeated.

\begin{table}[H]\centering\small
\begin{tabular}{p{3.6cm}p{2.7cm}p{6.1cm}p{1.8cm}}
\toprule
name & value & how estimated & split\\
\midrule
true wobble var.\ $\sigma_t^2$ & $0.00906\,\U^2$ & $\operatorname{cov}(r_A,r_C)$ of the trajectory-fit residuals, lines with $\ge5$ headline words, 21\,880 words, 1/99\% winsorising & train\\
headline noise var.\ $\sigma_A^2$ & $0.01176\,\U^2$ & $\operatorname{var}(r_A)-\sigma_t^2$ & train\\
mid-row noise var.\ $\sigma_C^2$ & $0.00631\,\U^2$ & $\operatorname{var}(r_C)-\sigma_t^2$ & train\\
headline weight $w$ & 0.349 & inverse-noise weight $\sigma_A^{-2}/(\sigma_A^{-2}+\sigma_C^{-2})$ & train\\
combined noise var. & $0.00411\,\U^2$ & $(\sigma_A^{-2}+\sigma_C^{-2})^{-1}$ & train\\
headline noise share & 0.565 (plausible 0.38--0.66) & $\sigma_A^2/\operatorname{var}(r_A)$; range from review M3 & train\\
combined noise share & 0.31 (0.29--0.33) & $\sigma^2_\text{comb}/(\sigma_t^2+\sigma^2_\text{comb})$ & train\\
wobble factor $f_\text{wobble}$ & $\approx0.83$ (0.82--0.84) & $\sqrt{1-\text{combined share}}$ & train\\
clearance $d_\text{min}$ & $2.603\,\U$ & 1st percentile of min.\ trajectory distance, 3\,740 consecutive line pairs (val: 2.613) & train\\
clearance $\rho_\text{min}$ & 0.797 & 1st percentile of min/mean distance (val: 0.788) & train\\
redraw limit & 30 & set by hand & ---\\
inflation$(n)$ (line scale) & 2.505 \ldots\ 1.234 ($n=3..12$) & simulation of i.i.d.\ offsets through \texttt{fit\_line}, 4000 lines per $n$, seed 0 & none\\
per-word factor & $\min(1/\sqrt{1-h_i},4)$ & OLS leverage of each word in the line's own fit; cap set by hand & none (per line)\\
per-word resid.\ var.\ $\bar s^2_n$ (curvature noise) & one value per $n$ & mean of $\text{RSS}/(n-3)$ over lines with $n$ words; $n\ge11$ pooled & train\\
curvature noise var.\ $v_i$ & $\bar s^2_{n_i}[(V_i^\top V_i)^{-1}]_{22}$ & OLS variance of $c_2$ from the line's own centred design & train ($\bar s^2_n$) / per line\\
curvature noise share & 0.384 & mean $v_i$ over var.\ of the measured curvature target, 4050 lines & train\\
rotation gain $\beta$ & 0.524 & Huber IRLS of headline tilt on tangent angle, 26\,653 words (val: 0.531) & train\\
\texttt{INK\_LEVEL} & 0.35 & set by hand to mimic the real binarisation (generated-page Otsu $\approx0.39$) & ---\\
input/target $\mu,\sigma$ & --- & standardisation statistics stored in the checkpoint & train\\
MDN weights $\theta$ (7 heads, $K=4$, width 64) & --- & maximum likelihood with teacher forcing; early stopping on val NLL (best epoch 152 of 183) & train (val for stopping)\\
\bottomrule
\end{tabular}
\caption{Learned and estimated parameters of the v2 default model
(\texttt{runs/default\_unfitw\_dc}). Sources: \src{runs/anchor\_noise.json},
\src{runs/line\_clearance.json}, \src{runs/unfit\_ratio.json},
\src{runs/default\_unfitw/tilt.json}, \src{runs/default\_unfitw/metrics.json},
\texttt{features.curv\_noise\_var} (curvature noise, \S\ref{sec:dc}),
review M3 (ranges). The deconvolved model has the same rotation gain and the
same best and stopping epochs as \texttt{default\_unfitw}. No constant was
fitted on test. Choosing \emph{between} designs, however, used the test split
(\S\ref{sec:limits}).}
\label{tab:params}
\end{table}

\begin{table}[H]\centering\small
\begin{tabular}{lccccccc}
\toprule
NLL (nats/target) & slope & curv & log res.\ scale & $x_\text{start}$ & pitch & gap & resid\\
\midrule
val  & $-2.363$ & $-0.975$ & 1.071 & 0.965 & 1.152 & 0.669 & $-0.606$\\
test & $-2.341$ & $-0.954$ & 1.073 & 0.844 & 1.176 & 0.679 & $-0.590$\\
\midrule
train rows & 4942 & 4050 & 4050 & 4814 & 4258 & 25\,546 & 26\,878\\
\bottomrule
\end{tabular}
\caption{Per-head NLL of the leverage model \texttt{default\_unfitw}
(\src{runs/default\_unfitw/metrics.json}). The default model
\texttt{default\_unfitw\_dc} is bit-identical on every head except curvature;
its validation curvature NLL, scored on the measured curvature under the
convolved likelihood~\eqref{eq:xd}, is $-1.004$ (\S\ref{sec:dc}). The NLLs
cannot be compared with those of models trained on different targets (fitted,
table or leverage residuals), because the targets themselves differ.}
\end{table}

%=====================================================================
\section{Evaluation protocol}
\label{sec:protocol}

\paragraph{Plan-level.} For each of the 193 test pages (unseen texts), every
generator receives the page's real word widths and line breaks. Its plans are
converted to raw coordinates and \emph{re-measured} by
\texttt{measure\_page}, the code that measured the real lines. The run is
repeated over 5 seeds, and we report mean $\pm$ sd across seeds. Each real page
and its generated copy form one group in the grouped 5-fold C2ST, so every
group holds both labels. Line C2ST features (lines with $n\ge5$): mean and sd
of the gaps, gap lag-1, slope, curvature, residual scale, drift, and the
correlation of gap with width. Page features: mean and sd of line-mean gaps,
mean and sd of slope, slope lag-1, mean curvature, mean residual scale, mean
pitch. The classifier is logistic regression unless the table says MLP.

\paragraph{Nulls.} The shuffled-real baseline (another test page's geometry)
scores 0.584 at line level. It is \emph{not} a real-vs-real floor: it pairs
donor gaps with the target page's widths, which breaks the gap--width relation
($0.018$ vs real $0.130$), and that relation is a C2ST feature (review M1). The
comparable within-page null splits each test page's lines into two labels in
the same group. It gives $0.496\pm0.019$ over 20 draws.

\paragraph{End-to-end.} Five reference writers (test pages image\_77, 104,
132, 147 and 910) $\times$ 2 seeds. BengaliScribe generates the words. The
layout is \texttt{manual} (BengaliScribe's constants), \texttt{v2}, or
\texttt{v2\_ctx} (conditioned on the writer's own real geometry). The rendered
pages go through the real detector and the same extraction, and are compared
with the 193 real test pages. From \texttt{e2e\_pagereset} onward the C2ST is
balanced (\S\ref{sec:balanced}). Recovered/placed is the ratio of extracted
words to drawn words.

\paragraph{2-D FFT (controlled).} A layout map is the filled word boxes at
8\,px/\U. It is windowed with a Hann window, and its 2-D power spectrum is
summarised as radial PSD (32 bins), angular energy (18 bins) and the vertical
line-pitch peak. In the \emph{controlled} version, every generator runs on each
real test page's own word widths. Real and generated pages are rendered with the
same fixed box (anchor $-0.25$ to $+0.85\,\U$) and cut to a $32\times18\,\U$
window inside the real 10th-percentile text block. Single-feature confound
AUCs are therefore 0.5 by construction.

\paragraph{1-D Fourier.} Dense headline traces of real lines and of
layout-realisable generated traces (rigid, rotated words) are compared by
log-spectral distance in the band with wavelength $\ge4\,\U$.

%=====================================================================
\section{Results}
\label{sec:results}

\subsection{Plan-level, 5 seeds}
\begin{table}[H]\centering
\resizebox{\linewidth}{!}{\begin{tabular}{lcccc}
\toprule
method & line C2ST & page C2ST & 2-D FFT C2ST (controlled) & resid.\ scale $W_1$, $n=5$ / 6--7 / 8--10 / 11+\\
\midrule
v1 (\texttt{scribe\_layout/runs/default}) & $0.566\pm0.021$ & $0.654\pm0.043$ & $0.567\pm0.034$ & .011 / .025 / .055 / .074\\
v2, old table (\texttt{v2\_history\_unfit}) & $0.539\pm0.016$ & $0.660\pm0.004$ & $0.653\pm0.024$ & .039 / .028 / .029 / .058\\
v2, leverage (\texttt{default\_unfitw}) & $0.521\pm0.019$ & $0.581\pm0.026$ & $0.566\pm0.025$ & .009 / .011 / .041 / .062\\
\textbf{v2, deconvolved (\texttt{default\_unfitw\_dc})} & $0.535\pm0.012$ & $0.590\pm0.024$ & $0.550\pm0.020$ & .007 / .011 / .042 / .061\\
shuffled real (not a floor) & $0.584\pm0.016$ & $0.605\pm0.012$ & $0.529\pm0.050$ & .021 / .018 / .054 / .100\\
\midrule
manual (BengaliScribe constants) & 0.980 & 1.0 & 0.99 & ---\\
within-page real-vs-real null & $0.496\pm0.019$ & --- & --- & ---\\
\bottomrule
\end{tabular}}
\caption{Plan-level comparison on the test split (193 pages, 1495 lines). The
v1 and old-table line and page C2ST values match \src{runs/eval/ablation.md}
(\texttt{learned\_full}, \texttt{v2\_history\_unfit}). The leverage row, the
FFT column and the residual-scale bins come from
\src{runs/compare\_unfitw/compare\_5seed.py} (5 seeds; output in
\src{runs/compare\_unfitw/compare\_5seed.txt}); the deconvolved row comes from
\src{runs/compare\_unfitw/compare\_5seed\_dc.py} (output in
\src{compare\_5seed\_dc.txt}). Bold: the current default. The FFT floor (real
halves) is 0.483 for logistic regression and 0.522 for the MLP.}
\label{tab:plan}
\end{table}

\paragraph{Measured.} The leverage model has the lowest line and page C2ST
among the learned models. Its line C2ST (0.521) is $0.025$ above the
within-page null (0.496), about 1.3 seed-sd. The deconvolved default model
scores $0.535\pm0.012$ (line), $0.590\pm0.024$ (page) and $0.550\pm0.020$
(FFT). Every difference from the leverage model is smaller than the larger
of the two seed-sds. The page C2ST of 0.58--0.59 still separates real from
generated pages better than chance. No page-level null was computed for this
protocol. The residual-scale $W_1$ is much lower than the old table's for
$n=5$--$7$, and still sizeable for long lines, for both leverage-family
models.

\begin{table}[H]\centering\small
\begin{tabular}{lcccc}
\toprule
$W_1$ to real (test, 5 seeds, mean) & slope & drift (\U) & curvature & pitch (\U)\\
\midrule
v2 leverage (\texttt{default\_unfitw}) & 0.0045 & 0.158 & 0.0002 & 0.109\\
v2 deconvolved (\texttt{default\_unfitw\_dc}) & 0.0044 & 0.155 & 0.0001 & 0.114\\
\bottomrule
\end{tabular}
\caption{Plan-level marginal $W_1$ of the two leverage-family models
(\src{runs/compare\_unfitw/drift.py}; output printed, not saved). Only the
curvature $W_1$ moves. The script also prints the seed-to-seed sd: 0.053 and
0.058 for drift, 0.029 and 0.039 for pitch, far larger than the
differences.}
\label{tab:drift}
\end{table}

\paragraph{Interpretation, with caveats.} The leverage correction removes the
per-word overshoot without giving up the line-level gain. The curvature
deconvolution leaves the plan-level C2ST and the slope, drift and pitch
marginals where they were and changes the line bow (\S\ref{sec:dc}). Line
C2ST 0.521--0.535 is \emph{close} to the null, but these models were chosen
after comparing alternatives on the same test split (\S\ref{sec:limits}). The
review estimated the selection bias at about 1 sd. On the validation split
(\S\ref{sec:val}) the same two models score 0.527 and 0.536, within seed noise
of their test values.

\paragraph{Earlier models (5 seeds, \src{runs/eval/ablation.md}).} Line/page
C2ST: \texttt{v2\_history} (fitted targets) $0.638/0.752$;
\texttt{v2\_latent} $0.650/0.749$; \texttt{v2\_latent\_unfit} $0.565/0.684$;
\texttt{v2\_history\_unfit} with writer context $0.540/0.574$. All learned
models reproduce the dependence structure that the manual layout lacks
(Fig.~\ref{fig:dep}): gap lag-1 about 0.47--0.49 (real 0.475), gap vs.\ width
(real 0.130) and slope persistence (real 0.884). The between-page ICC of
mean gap is too high in all of them (0.45--0.53 against real 0.312). So is the
clearance violation rate without rejection (0.09--0.15 against 0.014).

\begin{figure}[H]\centering
\includegraphics[width=\linewidth]{figs/dependence.png}
\caption{Dependence statistics on the test split, mean $\pm$ sd over 5 seeds
(\src{runs/eval/}; generated before the leverage model existed, so it is not
shown). Black: real. The manual layout (orange, third bar) has no gap
persistence, no gap--width relation and no slope persistence. The last panel is
the clearance-violation rate at plan level, without rejection sampling.}
\label{fig:dep}
\end{figure}

\subsection{Validation-split check}
\label{sec:val}
Every design choice above was compared on the test split (review M4). To see
how much of the test advantage survives on data that no comparison had used,
the plan-level protocol was re-run unchanged on the validation split (4
texts, 144 pages), with the same 5 seeds, the same generators and the same
controlled FFT set-up
(\src{runs/compare\_unfitw/compare\_5seed\_val.py}, output in
\src{compare\_5seed\_val.txt}). No design choice or fitted constant was
changed in response. The val split was not wholly unused before this check:
it sets the early-stopping epoch of the MDN weights, and the clearance
thresholds and the rotation gain were also computed on it as a cross-check
(Table~\ref{tab:params}).
\begin{table}[H]\centering
\resizebox{\linewidth}{!}{\begin{tabular}{lcccc}
\toprule
VAL split & line C2ST & page C2ST & FFT C2ST (contr.) & resid.\ scale $W_1$, $n=5$/6--7/8--10/11+\\
\midrule
v1 & $0.560\pm0.024$ & $0.631\pm0.040$ & $0.531\pm0.033$ & .013 / .025 / .055 / .089\\
v2, old table (\texttt{default\_unfit}) & $0.566\pm0.028$ & $0.660\pm0.022$ & $0.599\pm0.051$ & .049 / .030 / .024 / .050\\
v2, leverage (\texttt{default\_unfitw}) & $0.527\pm0.030$ & $0.575\pm0.039$ & $0.528\pm0.024$ & .012 / .011 / .035 / .072\\
\textbf{v2, deconvolved (\texttt{default\_unfitw\_dc})} & $0.536\pm0.027$ & $0.593\pm0.031$ & $0.535\pm0.044$ & .011 / .011 / .035 / .069\\
shuffled real (not a floor) & $0.577\pm0.023$ & $0.597\pm0.030$ & $0.480\pm0.034$ & .016 / .017 / .055 / .089\\
\bottomrule
\end{tabular}}
\caption{Plan-level comparison on the \emph{validation} split, 5 seeds,
mean $\pm$ sd. Same protocol as Table~\ref{tab:plan}. No within-page null was
computed on val.}
\label{tab:val}
\end{table}

\paragraph{Measured.} The two leverage-family models beat v1 on line C2ST
(0.527 and 0.536 against 0.560) and on page C2ST (0.575 and 0.593 against
0.631), as on test. On the controlled FFT C2ST they tie with v1 (0.528, 0.535,
0.531). The old-table model, which beat v1 on test line C2ST ($0.539$ against
$0.566$), does not on val ($0.566$ against $0.560$). Its FFT regression
(0.599) and its short-line residual-scale overshoot ($W_1$ 0.049 at $n=5$) do
replicate. The leverage model and the deconvolved model differ within seed
noise on every C2ST, as on test.

\paragraph{Interpretation, with caveats.} The improvement of the
leverage-family models over v1 on line and page C2ST is not an artefact of
selecting on test. The old table's test advantage over v1 on line C2ST
(\S\ref{sec:unfit}) was a selection effect, or at least did not generalise to
other texts, so it should not be cited as an improvement. The val margins over
v1 are about one seed-sd on line and page C2ST, and the split has only 4
texts, so the val numbers are noisy too. Having now been looked at, the val
split can no longer serve as an untouched split for future choices.

\subsection{FFT wobble-scale diagnostic}
\label{sec:fftdiag}
The old-table model improved line C2ST but \emph{worsened} the controlled 2-D
FFT C2ST, which suggested that its wobble was the wrong size. The diagnostic
scaled the drawn residuals of \texttt{v2\_history\_unfit} and re-ran the
controlled FFT C2ST (3 seeds, mean):
\begin{center}
\begin{tabular}{lcccccc}
\toprule
& v1 & v2 fitted targets & \multicolumn{4}{c}{v2 old table, drawn residual $\times$}\\
& & (\texttt{runs/default}) & 1.0 & 0.83 & 0.5 & 0.0\\
\midrule
2-D FFT C2ST & 0.612 & 0.628 & 0.683 & 0.658 & 0.578 & 0.601\\
\bottomrule
\end{tabular}
\end{center}
\paragraph{Measured.} The AUC falls as the drawn wobble is scaled down to
$\times0.5$, then rises again when the wobble is removed entirely.
\paragraph{Interpretation.} The page spectrum detects the excess per-word wobble
of the old table. Some wobble is still needed, since $\times0$ is worse than
$\times0.5$. This motivated the leverage correction of \S\ref{sec:lev}, which
lands at $0.566\pm0.025$ in the 5-seed run. These are 3-seed numbers from a
separate script (\src{runs/compare\_unfitw/fft\_wobble\_scale.py}, output in \src{fft\_wobble\_scale.txt}), and single-seed FFT AUCs move by up to
$\pm0.09$. Compare them only with each other, not with Table~\ref{tab:plan}.

\subsection{2-D FFT layout-map experiment}
\label{sec:fft2d}
\paragraph{Uncontrolled version: confounded.} Detector-extracted real pages were
compared with BengaliScribe-rendered pages. The spectral C2ST AUCs (logistic) were
0.958--1.0 for every method (\src{scribe\_layout\_fft2d/runs/results.md}).
However, single non-layout features separate the sources just as well: text
block height (single-feature AUC 0.008--0.017, i.e.\ 0.98--0.99 after flipping
the direction, for manual, v1 and v2), median word-box width (0.90--0.93 after
flipping, for v1, v2 and manual), and box height (about 0.81--0.82 for
\texttt{v2} and \texttt{v2\_old}). The uncontrolled AUCs measure text,
page fill and glyph shape, not layout, and we draw no layout conclusion from
them.

\paragraph{Controlled version: the main result.} With text-matched word widths,
a fixed box convention and a common window, all confound AUCs are 0.5 for
every method. On a single seed (\src{results.md}), logistic C2ST was: manual
0.992, shuffled real 0.552, v1 0.552, \texttt{v2\_history\_unfit} 0.661,
\texttt{v2\_latent\_unfit} 0.685. The real-halves floor was 0.483. The 5-seed
values in Table~\ref{tab:plan} supersede these: v1 $0.567\pm0.034$, old table
$0.653\pm0.024$, leverage $0.566\pm0.025$, shuffled real $0.529\pm0.050$.
Fig.~\ref{fig:fft} shows that the mean log-PSD of every learned model lies on
top of the real one, while the manual layout has a strong spurious peak
structure at low frequency (its regular line pitch).

\paragraph{What the spectrum is sensitive to.} We ablated real pages, pooling
the population. Shuffling the line pitch moves the mean log-PSD 6 times more
(L2 0.080) than zeroing the slope (0.013) or removing the wobble (0.013). The
page spectrum is therefore dominated by line-pitch regularity. The
peak-prominence statistic is unreliable: it \emph{rises} after periodicity is
destroyed ($243\to273$).

\begin{figure}[H]\centering
\includegraphics[width=0.55\linewidth]{figs/mean_log_psd_fair.png}\\[4pt]
\includegraphics[width=\linewidth,trim=0 60 0 60,clip]{figs/example_layout_maps_fair.png}
\caption{Controlled 2-D FFT experiment
(\src{scribe\_layout\_fft2d/runs/figures/}). Top: mean log radial PSD per
source, 193 pages each. All learned sources and shuffled real lie under the
real curve, and manual (orange) departs from it at low frequency. Bottom: one
page's layout map and log power spectrum per source. Learned models are
included up to the old-table model, before the leverage model existed.}
\label{fig:fft}
\end{figure}

\subsection{1-D Fourier experiment on headline traces}
\label{sec:fourier1d}
\paragraph{Method.} Dense headline traces (\texttt{scribe\_layout\_fourier/})
were analysed by Lomb--Scargle PSD. Real wobble follows $\sim1/f^\alpha$ with
$\alpha=1.06$ ($R^2=0.99$, train, 5099 lines). A spectral-synthesis generator
draws traces with this spectrum.

\paragraph{Fair comparison.} Words are rigid, so a layout can only realise a
trace of rotated, rigid word segments. We compared layout-realisable traces in
the band of wavelength $\ge4\,\U$ by log-spectral distance to real
(\src{runs/fair\_dense\_psd.json}):
\begin{center}
\begin{tabular}{lcccc}
\toprule
& manual & v1 & Fourier & Fourier, continuous (not realisable)\\
\midrule
distance to real & 3.786 & 0.702 & 0.668 & 0.244\\
\bottomrule
\end{tabular}
\end{center}
\paragraph{Interpretation.} Once the comparison is restricted to traces a
word-placement layout can actually produce, the Fourier generator and v1 are
close (0.668 vs 0.702). The large advantage seen earlier holds only for
continuous traces that rigid words cannot realise. An earlier claim that the
Fourier model was ``17$\times$ better'' (3.98 vs 0.23 nats,
\src{scribe\_layout\_fourier/runs/results.md}) was an artefact of comparing
against that unrealisable continuous trace. A claimed ``6$\times$ better
$x_\text{start}$'' came from including paragraph indents. At word level the
Fourier model is worse on gaps ($W_1$ 0.092 vs 0.038) and cannot reproduce the
gap--width or gap--position relations. This matches the survey's prediction
that spectra suit dense traces, not 4--8 gaps.

\subsection{End-to-end}
\label{sec:e2e}
Tables~\ref{tab:e2e-biased} and~\ref{tab:e2e-balanced} separate the runs that used the biased pooled C2ST from those that used the balanced one. Do not compare C2ST values across the two tables.
\begin{table}[H]\centering\small
\begin{tabular}{llcccc}
\toprule
run (placement) & method & gap $W_1$ & line C2ST & page C2ST & recovered\\
\midrule
\texttt{e2e\_unfit} & manual & 0.685 & 0.893 & 0.998 & 1.004\\
(canvas width) & v1     & 0.582 & 0.699 & 0.942 & 0.797\\
 & v2     & 0.602 & 0.750 & 0.901 & 0.969\\
 & v1\_ctx & 0.526 & 0.653 & 0.971 & 0.714\\
 & v2\_ctx & 0.381 & 0.595 & 0.901 & 0.910\\
\midrule
\texttt{e2e\_unfit\_ink} & v2     & 0.471 & 0.667 & 0.840 & 0.952\\
(ink extents) & v2\_ctx & 0.207 & 0.532 & 0.918 & 0.866\\
\midrule
\multicolumn{2}{l}{real-vs-real ``floor'' (biased)} & --- & 0.341 & 0.115 & ---\\
\bottomrule
\end{tabular}
\caption{\textbf{Biased-C2ST runs.} End-to-end runs with the pooled grouped C2ST, which is biased low
(\src{runs/e2e\_unfit/results.md}, \src{runs/e2e\_unfit\_ink/results.md}).
Model: \texttt{default\_unfit} (old table), with no planner reset. The
real-vs-real values far below 0.5 show that the protocol was broken. The
$W_1$ and recovered columns are not affected by the C2ST bias. They do depend
on which lines survive extraction.}
\label{tab:e2e-biased}
\end{table}

\begin{table}[H]\centering
\resizebox{\linewidth}{!}{\begin{tabular}{llcccccccc}
\toprule
run (model) & method & gap $W_1$ & slope $W_1$ & drift $W_1$ & pitch $W_1$ & $x_\text{start}$ $W_1$ & line C2ST & page C2ST & recovered\\
\midrule
\texttt{e2e\_pagereset} & manual & 0.685 & 0.035 & 0.965 & 1.839 & 0.948 & 0.905 & 0.999 & 1.004\\
(\texttt{default\_unfit}) & v2     & 0.404 & 0.013 & 0.396 & 0.670 & 4.065 & 0.684 & 0.834 & 0.965\\
 & v2\_ctx & 0.267 & 0.017 & 0.558 & 0.770 & 3.783 & 0.619 & 0.889 & 0.926\\
\midrule
\texttt{e2e\_unfitw} & manual & 0.685 & 0.035 & 0.965 & 1.839 & 0.948 & 0.905 & 0.999 & 1.004\\
(\texttt{default\_unfitw}) & v2     & 0.376 & 0.012 & 0.343 & 0.346 & 2.278 & 0.699 & 0.851 & 0.960\\
 & v2\_ctx & 0.197 & 0.019 & 0.610 & 0.480 & 2.703 & 0.613 & 0.842 & 0.933\\
\midrule
\texttt{e2e\_unfitw\_dc} & manual & 0.685 & 0.035 & 0.965 & 1.839 & 0.948 & 0.905 & 0.999 & 1.004\\
(\texttt{default\_unfitw\_dc}) & v2     & 0.387 & 0.021$^\dagger$ & 0.645$^\dagger$ & 0.465 & 1.453 & 0.719 & 0.844 & 0.983\\
 & v2\_ctx & 0.183 & 0.023$^\dagger$ & 0.695$^\dagger$ & 0.535 & 2.494 & 0.654 & 0.847 & 0.929\\
\midrule
\multicolumn{2}{l}{real-vs-real floor (balanced, 5 seeds)} & --- & --- & --- & --- & --- & 0.503 & 0.527 & ---\\
\bottomrule
\end{tabular}}
\caption{\textbf{Balanced-C2ST runs.} End-to-end runs with the balanced C2ST, planner reset per page and
ink-extent placement. The three runs use the same protocol, writers and seeds;
only the layout model differs. Sources: \src{runs/e2e\_pagereset/results.md}
(run floor 0.534 line, 0.508 page), \src{runs/e2e\_unfitw/results.md} and
\src{runs/e2e\_unfitw\_dc/results.md} (same run floor, 0.534 and 0.508). The
manual rows are identical because the manual layout does not use the model.
The line C2ST features omit $x_\text{start}$ and pitch. $^\dagger$All $W_1$
values are computed on the lines that survive extraction. Slope and drift
$W_1$ in particular are not comparable across runs with different split rates
(see the survivor-bias analysis below). The \texttt{e2e\_unfitw} slope and
drift values were flattered by excluding split lines; the
\texttt{e2e\_pagereset} values very probably were too, but its split lines
were not analysed.}
\label{tab:e2e-balanced}
\end{table}

\paragraph{Leverage model end to end (\texttt{e2e\_unfitw}).} Compared with
\texttt{e2e\_pagereset}, where only the model differs, the leverage model has
lower gap $W_1$ (v2 $0.404\to0.376$; v2\_ctx $0.267\to0.197$), lower pitch
$W_1$ ($0.670\to0.346$; $0.770\to0.480$) and lower $x_\text{start}$ $W_1$
($4.065\to2.278$; $3.783\to2.703$). Drift $W_1$ is better for v2
($0.396\to0.343$) and worse for v2\_ctx ($0.558\to0.610$). Line C2ST is
essentially unchanged (v2 $0.684\to0.699$; v2\_ctx $0.619\to0.613$). Page
C2ST moves in opposite directions (v2 $0.834\to0.851$; v2\_ctx
$0.889\to0.842$) and stays far above the floor. The extractor found fewer lines
on the generated pages (v2 $499\to456$, v2\_ctx $461\to454$) for about the
same number of placed words. The pitch and $x_\text{start}$ gains are
\emph{consistent with} less line fragmentation. Drawn-versus-extracted line
counts were later measured for \texttt{e2e\_unfitw} (below) but not for
\texttt{e2e\_pagereset}, so this comparison remains an interpretation.
With 10 generated documents per method, changes in C2ST of a few hundredths
are within noise.

\paragraph{Line splitting against the drawn lines.} Matching extracted lines
to the drawn lines (\S\ref{sec:dc}; \src{runs/compare\_unfitw/split\_diag.py},
method v2):
\begin{center}\small
\begin{tabular}{lccccc}
\toprule
run (v2) & drawn lines & drawn lines split & single-piece extr.\ lines & extra fragments\\
\midrule
\texttt{e2e\_unfitw} & 462 & 89 (19\%) & 127 & 113\\
\texttt{e2e\_unfitw\_dc} & 469 & 52 (11\%) & 64 & 61\\
\texttt{e2e\_final} (guard) & 453 & 57 (12.6\%) & 80 & 69\\
\bottomrule
\end{tabular}
\end{center}
The manual layout in \texttt{e2e\_unfitw} has 6 of 355 drawn lines split
(1.7\%). The deconvolved model halves the single-piece lines and the extra
fragments, and the v2 extraction keeps more lines (473 against 456) with a
higher recovered/placed ratio (0.983 against 0.960).

\paragraph{Deconvolved default model end to end (\texttt{e2e\_unfitw\_dc}).}
Compared with \texttt{e2e\_unfitw}, $x_\text{start}$ $W_1$ falls (v2
$2.278\to1.453$; v2\_ctx $2.703\to2.494$), consistent with fewer mid-page
fragments. Gap $W_1$ moves little ($0.376\to0.387$; $0.197\to0.183$). Line
C2ST rises slightly (v2 $0.699\to0.719$; v2\_ctx $0.613\to0.654$) and page
C2ST is unchanged (0.844 and 0.847). Slope $W_1$ ($0.012\to0.021$;
$0.019\to0.023$), drift $W_1$ ($0.343\to0.645$; $0.610\to0.695$) and pitch
$W_1$ ($0.346\to0.465$; $0.480\to0.535$) rise. At plan level, where nothing
is split, the two models have the same slope, drift and pitch $W_1$
(Table~\ref{tab:drift}).

\paragraph{Survivor bias in the end-to-end slope and drift.} The drawn slope
of every generated line can be read from its drawn line polygon
(\src{runs/compare\_unfitw/surv.py}, method v2; top edge of the polygon, with
the bow as the range of the detrended top edge). In \texttt{e2e\_unfitw}, the
lines that the extraction split have median $|$slope$|$ 0.046 (90th percentile
0.096) and median bow $0.66\,\U$; the lines it kept whole have 0.026 (0.070)
and $0.32\,\U$. In \texttt{e2e\_unfitw\_dc}, split lines have median
$|$slope$|$ 0.052 and kept lines 0.025. Over \emph{all} drawn lines of
\texttt{e2e\_unfitw\_dc} the median $|$slope$|$ is 0.027 (90th percentile
0.071), against 0.029 (0.086) measured on real test lines with $n\ge3$.

\emph{Measured:} split lines are the steep and bowed ones, and the drawn
slopes as a whole are close to the real ones. \emph{Interpretation:} the
end-to-end slope and drift $W_1$ are computed only on the surviving lines.
When 19\% of the lines are split, the steep and bowed lines drop out of the
metrics, the survivors look straighter and flatter than the generator's
output, and the $W_1$ to real is flattered. When fewer lines are split, more
of the steep lines enter the metrics. The rise in slope and drift $W_1$ from
\texttt{e2e\_unfitw} to \texttt{e2e\_unfitw\_dc} is therefore not evidence
that the deconvolved model draws worse slopes, and the earlier end-to-end
slope and drift numbers (Table~\ref{tab:e2e-balanced}) were flattered by
excluding split lines. End-to-end slope and drift $W_1$ are not comparable
across runs with different split rates. The drawn and real slopes come from
different measurements (polygon top edge vs.\ trajectory fit), so their
closeness is approximate. The rise in pitch $W_1$ is not explained by this
analysis.

\paragraph{Final default end to end (\texttt{e2e\_final}).}
Table~\ref{tab:e2e-final} is the end-to-end result of the final default
configuration: model \texttt{runs/default\_unfitw\_dc} with the margin guard
of \S\ref{sec:margin}, same writers, seeds, balanced C2ST and real reference
as Table~\ref{tab:e2e-balanced}. The layout code differs from
\texttt{e2e\_unfitw\_dc} only by the margin guard.
\begin{table}[H]\centering
\resizebox{\linewidth}{!}{\begin{tabular}{lccccccccccc}
\toprule
method & gap $W_1$ & slope $W_1$ & drift $W_1$ & pitch $W_1$ & $x_\text{start}$ $W_1$ & line C2ST & page C2ST & recovered & pages & lines kept\\
\midrule
manual & 0.685 & 0.035 & 0.965 & 1.839 & 0.948 & 0.905 & 0.999 & 1.004 & 20 & 327\\
\textbf{v2}     & 0.425 & 0.021 & 0.672 & 0.390 & 1.239 & 0.719 & 0.850 & 0.981 & 30 & 440\\
\textbf{v2\_ctx} & 0.198 & 0.020 & 0.621 & 0.503 & 2.130 & 0.619 & 0.845 & 0.940 & 33 & 432\\
\midrule
run real-vs-real floor & --- & --- & --- & --- & --- & 0.534 & 0.508 & --- & --- & ---\\
\bottomrule
\end{tabular}}
\caption{\textbf{Final default end-to-end result} (\src{runs/e2e\_final/results.md}):
default model \texttt{default\_unfitw\_dc} plus the margin guard; balanced
C2ST; 5 writers $\times$ 2 seeds; real reference 193 test pages. All $W_1$
values are on the lines that survive extraction (survivor bias, see above).
Against the drawn lines (\src{split\_diag.py}, v2): 57 of 453 drawn lines
split (12.6\%), 80 single-piece extracted lines, 69 extra fragments.}
\label{tab:e2e-final}
\end{table}
\emph{Measured:} against \texttt{e2e\_unfitw\_dc}, line C2ST is 0.719 for v2
in both runs and moves $0.654\to0.619$ for v2\_ctx; page C2ST moves
$0.844\to0.850$ and $0.847\to0.845$. The split rate of drawn v2 lines is
12.6\% (57 of 453) against 11.1\% (52 of 469). The $W_1$ values move in both
directions: gap $0.387\to0.425$ and $0.183\to0.198$, drift $0.645\to0.672$
and $0.695\to0.621$, pitch $0.465\to0.390$ and $0.535\to0.503$,
$x_\text{start}$ $1.453\to1.239$ and $2.494\to2.130$. The v2 method placed
2249 words against 2339. With 10 generated documents per method and no
end-to-end seed-sd, we read the two runs as the same within noise.
\emph{Interpretation:} a replanned line very probably changes the random
stream for everything drawn after it, so the two runs are best seen as
different draws of nearly the same layout process, not as a paired
comparison of the guard. The guard's effect is shown by the margin test
(\S\ref{sec:margin}), not by these numbers.


\paragraph{Measured.} With the balanced test, line C2ST ranges from 0.613 to
0.719 for the learned layouts and is 0.905 for manual. Page C2ST is 0.83--0.89
for the learned layouts, so real and generated pages remain easy to tell apart.
The learned layouts reduce gap, slope, drift and pitch $W_1$ well below manual
(on the surviving lines). $x_\text{start}$ $W_1$ is worse than manual:
3.8--4.1\,\U\ with the old-table model, 2.3--2.7\,\U\ with the leverage
model and 1.2--2.5\,\U\ with the deconvolved model (with and without the
margin guard), against 0.95\,\U.

\paragraph{Interpretation, with caveats.} The large $x_\text{start}$ $W_1$ comes
mostly from fragments of generated lines that the extractor splits apart: a
fragment starts mid-page. Part of it is real margin spread. BengaliScribe also
adds paragraph gaps and indents that the real hcr pages lack (review minor 7).
The extraction splits generated lines heavily. In the pre-reset run
\texttt{e2e\_unfit\_ink}, v2\_ctx drew 507 lines; the extractor found 829 and
kept 470. 39\% of extracted lines were single pieces, against 12.9\% on real
pages. Against the drawn lines, 19\% of v2 lines were split with the leverage
model and 11\% with the deconvolved model (12.6\% in \texttt{e2e\_final}),
and the split lines are the steep
or bowed ones. The metrics therefore see only the lines that survive
extraction, and that subset is biased. Only
about 50--60 generated lines with $n\ge5$ per seed enter the line C2ST.
Fig.~\ref{fig:e2e} (middle) shows a v2\_ctx page from the old-table model,
where some words jump visibly off the line. This is the per-word overshoot of
\S\ref{sec:unfit}, which the leverage model is designed to remove.

\begin{figure}[H]\centering
\includegraphics[width=0.32\linewidth]{figs/pagereset_image_104_manual.png}\hfill
\includegraphics[width=0.32\linewidth]{figs/pagereset_image_104_v2_ctx.png}\hfill
\includegraphics[width=0.32\linewidth]{figs/unfitw_image_104_v2_ctx.png}
\caption{End-to-end pages for reference writer image\_104
(\src{runs/e2e\_pagereset/thumbs/}, \src{runs/e2e\_unfitw/thumbs/}).
Overlays are the lines re-extracted by the real detector and extraction. Left:
manual layout, with straight, parallel lines and uniform spacing. Middle: v2
with writer context, model \texttt{default\_unfit} (old table). Note the
per-word vertical jumps, for example the raised last two words of line 2.
Right: the same words, seed and context with the leverage model
\texttt{default\_unfitw}. The jump in line 2 is smaller, while the raised
pair of words in the ninth line remains. One page cannot show a population
effect. The figure illustrates the change and is not evidence for it.}
\label{fig:e2e}
\end{figure}

\subsection{Qualitative BengaliScribe comparisons}
\begin{figure}[H]\centering
\includegraphics[width=\linewidth]{figs/compare_v1_v2.png}
\caption{The same generated words under four layouts, top of page 1, reference
writer image\_132 (\src{runs/bscribe/compare\_v1\_v2.png}). From left: manual,
v1, v2 (old-table model \texttt{default\_unfit}), v2 with writer context.
Rendered before the planner reset and before the leverage model.}
\label{fig:v1v2}
\end{figure}

\begin{figure}[H]\centering
\includegraphics[width=\linewidth]{figs/compare_fixes.png}
\caption{Re-layout of the same generated words (\texttt{relayout.py};
\src{runs/bscribe/compare\_fixes.png}). Left: base (no de-noising, no
clearance). Middle: wobble de-noised ($\times0.83$). Right: de-noising plus
line clearance. In the right panel the lines no longer converge or touch. The
de-noising effect is subtle at this scale. Rendered with the model before the
unfit change.}
\label{fig:fixes}
\end{figure}

%=====================================================================
\section{Corrections from independent review}
\label{sec:review}
An independent read-only review (\texttt{runs/review\_opus.md}) reproduced each
claim with a script (\texttt{c1}--\texttt{c19.py}). The table lists what it
found and what changed. VERIFIED means reproduced by script. The last two
rows are not from the review: they were found by the tests written in
response to its ``tests'' finding.

{\small
\setlength{\LTcapwidth}{\linewidth}
\begin{longtable}{>{\raggedright\arraybackslash}p{1.1cm}>{\raggedright\arraybackslash}p{5.0cm}>{\raggedright\arraybackslash}p{3.4cm}>{\raggedright\arraybackslash}p{5.0cm}}
\toprule
id & what was wrong & how it was found & what changed\\
\midrule
\endfirsthead
\toprule
id & what was wrong & how it was found & what changed\\
\midrule
\endhead
\midrule
\multicolumn{4}{r}{\footnotesize continued on the next page}\\
\endfoot
\bottomrule
\caption{Review findings and their status. Minor findings 1, 4, 5, 6, 8 and 9
are listed in \texttt{review\_opus.md}. Minor 1 found the de-noising factor
does not double-count. Minors 4--5 are training/sampling conditioning
mismatches for short lines, suspected but not shown to matter. Minor 6 is
unfiltered context lines. Minor 8 is writer overlap. Minor 9 is the inclusion
of 5 context pages in the real reference.}\\
\endlastfoot
C1 & End-to-end C2ST biased low by single-label groups; the reported v2\_ctx line C2ST 0.532 did not hold (0.627 with per-fold AUC) & real-vs-real null 0.457$\pm$0.067 pooled vs 0.506$\pm$0.065 per-fold (20 draws); fold-seed range 0.41--0.61 & balanced folds with per-fold mean AUC (\S\ref{sec:balanced}); biased and balanced runs reported separately\\
C2 & Extraction fragments generated lines; e2e metrics see only the surviving lines & line counts drawn, extracted and kept; re-render with zero wobble still fragmented & \emph{reduced, not fixed}: slope drift (C3) and over-bowed lines (\S\ref{sec:dc}) were causes; 19\%\,$\to$\,11\% of drawn v2 lines split, 12.6\% in \texttt{e2e\_final} (\S\ref{sec:e2e})\\
C3 & Planner history ran across pages; slopes random-walked to the clamp; words drawn past the page bottom & tracing \texttt{hist.count} (56--72) and per-page slope means & reset per page to a fixed reference state (\S\ref{sec:reset})\\
M1 & Shuffled real (0.584) presented as a real-vs-real floor & gap--width correlation 0.018 vs 0.130; within-page null & null reported as $0.496\pm0.019$; ``below floor'' wording dropped\\
M2 & Unfit table applied to per-word targets: per-word RMS $2.07\times$ physical at $n=5$ & simulation; resid lag-1 and length dependence of real wobble & per-word leverage correction (\S\ref{sec:lev}), on which the default model is built; serial correlation still not modelled\\
M3 & ``57\% noise'' rests on independence of estimator errors, which fails & tetrad test with 2 extra estimators; clip-level sweep & headline share reported as 38--66\%; the combined share 0.29--0.33 used for generation\\
M4 & All design choices compared on test; val never used & reading \texttt{evaluate.py} & \emph{partly addressed}: an after-the-fact check on val (\S\ref{sec:val}); designs still chosen on test (\S\ref{sec:limits})\\
minor 2 & \texttt{INK\_LEVEL} 0.35 vs generated-page Otsu $\approx0.39$: gaps $+0.067\,\U$ & measuring ink edges & not fixed; documented\\
minor 3 & Clearance check ran before the page-break decision & code reading & current code redraws the line shape under the new page's state when a line moves to a new page\\
minor 7 & BengaliScribe adds paragraph gaps and indents that real pages lack & pitch after short lines on real pages (median 3.87 vs 4.44) & not changed; caveat on pitch and $x_\text{start}$\\
minor 10 & \texttt{anchor\_noise.json} stale against the cache (21\,880 vs 21\,794 words; $w$ 0.349 vs 0.351); circular dependency & recomputation & negligible effect; reproduction order matters\\
tests & Tests that pass vacuously: noise-floor tolerance $|\text{AUC}-0.5|<0.4$; one-sample groups in the C2ST test; latent test on a fake model; no tests for \texttt{bscribe\_layout}, \texttt{unfit}, \texttt{anchor\_noise}, \texttt{line\_clearance} & reading the tests against pathological outputs & \emph{fixed or explained}: 70 tests, about 17\,s (was 32, about 8\,min). Noise-floor tolerance $0.4\to0.15$ (measured max.\ deviation about 0.08). Balanced grouped-$k$-fold tests with multi-row single-label groups (real-vs-real near 0.5; \texttt{balanced=False} biased low; shifted data $>0.8$). New \texttt{test\_unfit}, \texttt{test\_bscribe\_layout} (placement order across pages, margins, exact per-page history reset with and without context, clearance stats), \texttt{test\_line\_clearance}, \texttt{test\_anchor\_noise} (re-implements \texttt{main()}'s formulas on synthetic data, since \texttt{main} needs the real data). The latent test is kept on its fake model, with a docstring saying what it does not show\\
tests (new) & Margin bug: a sloped first line rose 1--1.4\,\U\ into the top margin; ink could pass the bottom margin & \texttt{TestWordsStay\allowbreak InsideMargins}, which fails on the old code & \emph{fixed}: margin guard (\S\ref{sec:margin})\\
tests (new) & Leverage factor under-corrects edge words by about 10--25\% under the Huber fit & \texttt{test\_unfit.py} & tried a Huber-aware factor; worse on every C2ST; \emph{not adopted}, open (\S\ref{sec:huber})\\
\end{longtable}
}

The review's verdicts: plan-level line C2ST survives with caveats (correct null
about 0.50; selected on test). The end-to-end line C2ST of 0.532 did not hold.
The end-to-end gap $W_1$ survives as a measured number on the surviving lines,
with a $+0.07\,\U$ ink-edge bias. The noise share needs a range.

%=====================================================================
\section{Limitations and threats to validity}
\label{sec:limits}
\begin{itemize}
\item \textbf{Selection on test.} \texttt{evaluate.py} evaluates only on the
  test split. Fifteen methods were compared on test, and the unfit, de-noise,
  ink-extent, clearance and leverage choices were all iterated against test
  numbers. The fitted constants come from train (Table~\ref{tab:params}), so
  nothing leaked into the fits, but the test numbers are development-set
  numbers. The validation check (\S\ref{sec:val}) softens this: with the
  designs frozen, the leverage-family models score 0.527 and 0.536 on val
  line C2ST against v1's 0.560, close to their test values, and beat v1 on
  page C2ST too. It also shows the risk was real: the old table's test
  advantage over v1 did not replicate on val. The val split informed no design
  choice, but it was looked at after the fact, it sets the early-stopping
  epoch, and it has only 4 texts. It is therefore a replication, not a clean
  held-out estimate, and it cannot serve as an untouched split for further
  choices. The Huber-aware correction (\S\ref{sec:huber}) was then judged
  on both test and val; it was rejected, so this does not flatter the
  default, but no untouched split remains.
\item \textbf{No writer IDs.} There are 996 pages, about 37 copies of each of
  27 texts, scanned in 57 batches, with no writer identifiers. The same people
  very probably wrote several texts, so the text-level split probably does not
  separate writers. Test numbers may therefore include writer overlap with
  train. This cannot be checked with the available data.
\item \textbf{End-to-end line splitting.} The extractor splits generated
  lines far more often than real ones. Against the drawn lines, 11\% of the
  default model's v2 lines are split (52 of 469), down from 19\% (89 of 462)
  with the leverage model, and 12.6\% (57 of 453) in the final run with the
  margin guard; the manual layout's rate is 1.7\%. With writer context the
  rate is higher: 20\% (92 of 464) in the final run, down from 26\% (125 of
  478) with the leverage model; some context writers slope close to the
  extractor's 0.15 slope limit (e.g.\ image\_77, mean slope about $-0.15$
  on generated pages). The split lines are the steep or bowed ones
  (\S\ref{sec:e2e}); counts for all runs and variants are in
  \src{runs/compare\_unfitw/split\_counts.txt}. End-to-end metrics are
  computed on the surviving lines only, so slope and drift $W_1$ are biased
  towards straight lines and are not comparable across runs with different
  split rates, and $x_\text{start}$ $W_1$ is inflated by fragments. The
  planner reset and the curvature deconvolution address two causes. The lines
  that still split are again the steeper ones (median $|$slope$|$ 0.052), but
  why the chainer fails on them has not been analysed for this run.
\item \textbf{The latent model regression is unresolved.} The causal-prefix
  fix helps context generation but hurts unconditional generation
  (page C2ST $0.684\to0.829$). The latent model is not used.
\item \textbf{Noise-share range.} The headline-only noise share lies between
  38\% and 66\%, depending on assumptions that the data partly reject. The
  generation-time factor (0.82--0.84) is robust to this range. Any statement
  about how much measured wobble is ``real'' is not.
\item \textbf{Small end-to-end samples.} Five writers $\times$ 2 seeds give
  20--35 generated pages per method and only about 50--60 lines with $n\ge5$ per
  seed in the line C2ST. A single fold seed moved one old end-to-end line AUC
  between 0.41 and 0.61. End-to-end differences of a few hundredths are
  within noise.
\item \textbf{Wobble model (open).} Per-word wobble is drawn i.i.d., but real
  wobble is smoother (residual lag-1 $-0.38$ against about $-0.48$ for
  i.i.d.) and grows with line length. Long lines remain off ($W_1$ 0.061 at
  $n\ge11$ for the default). The per-word correction is calibrated for i.i.d.\
  offsets under least squares; it under-corrects edge words under the actual
  Huber fit, and the exact Huber-aware i.i.d.\ correction scores worse
  (\S\ref{sec:huber}), which suggests the i.i.d.\ assumption is the larger
  error. Calibrating the correction under a correlated-wobble model fitted to
  the real lag correlations is the open fix; it has not been tried.
\item \textbf{Margins (fixed).} Before the margin guard, a sloped or bowed
  first line could rise about 1--1.4\,\U\ into the top margin, and ink could
  pass the bottom margin. The guard (\S\ref{sec:margin}) fixes this, and a
  test that fails on the old code now covers it. It reserves only
  $0.3\,\U$ for the wobble, so an extreme word can still cross a margin by a
  little; the test allows $0.5\,\U$.
\item \textbf{Tests.} The 70 unit tests (\S\ref{sec:review}) check
  code paths and invariants on synthetic data. They do not check that a
  trained model reproduces real statistics; that is what the plan-level and
  end-to-end evaluations are for. The latent mean-reversion test still uses a
  fake model.
\item \textbf{Comparability across runs.} The leverage, deconvolved, Huber-aware and
  validation plan-level numbers come from comparison scripts, not from the
  standard \texttt{evaluate.py} run directory. Some of these scripts
  (\texttt{sag.py}, \texttt{drift.py}, \texttt{curvnoise2.py},
  \texttt{split\_diag.py}, \texttt{split\_why.py}, \texttt{surv.py}) print
  their output without saving it. They use the same harness functions
  (\texttt{evaluate.generate}, \texttt{evaluate.c2st}), but the metrics files
  in \texttt{runs/eval} do not contain them.
\end{itemize}

%=====================================================================
\section{Reproduction and files}
From \texttt{/home/hpc4090/shyan/ocr}:
\begin{lstlisting}
venv/bin/python -m scribe_layout_v2.anchor_noise           # runs/anchor_noise.json (train)
venv/bin/python -m scribe_layout_v2.extract                # re-extract with combined anchor
venv/bin/python -m scribe_layout_v2.line_clearance         # runs/line_clearance.json
venv/bin/python -m scribe_layout_v2.unfit                  # runs/unfit_ratio.json
venv/bin/python -m scribe_layout_v2.train --config scribe_layout_v2/configs/default_unfitw.yaml
venv/bin/python -m scribe_layout_v2.train --config scribe_layout_v2/configs/default_unfitw_dc.yaml  # default
venv/bin/python -m scribe_layout_v2.fit_tilt               # runs/<name>/tilt.json
venv/bin/python -m scribe_layout_v2.evaluate               # plan-level, 5 seeds -> runs/eval
venv/bin/python -m scribe_layout_v2.e2e_check --seeds 2 \
    --model scribe_layout_v2/runs/default_unfitw_dc/model.pt \
    --out scribe_layout_v2/runs/e2e_final --variants manual v2 v2_ctx   # end-to-end (final)
venv/bin/python -m pytest -q scribe_layout_v2/tests        # 70 tests, about 17 s
venv/bin/python -m scribe_layout_v2.bscribe_run --ref hcr/996/Images/image_132.png \
    --out scribe_layout_v2/runs/bscribe/v2_final --layout learned --context image_132
\end{lstlisting}
\texttt{anchor\_noise} reads the extraction cache and \texttt{extract} reads
\texttt{anchor\_noise.json} (review minor 10), so the order matters for exact
reproduction. \texttt{bscribe\_run.py} and \texttt{relayout.py} default to
\texttt{runs/default\_unfitw\_dc/model.pt}. The rejected Huber-aware model
is trained with \texttt{--config scribe\_layout\_v2/configs/default\_unfitwh\_dc.yaml};
\texttt{python -m scribe\_layout\_v2.unfit} also writes its factor table
\texttt{runs/unfit\_word\_huber.json}. The test suite (70 tests, about
17\,s on this machine; 32 tests and about 8\,min before) runs with the
pytest command above.

The comparison scripts in \texttt{runs/compare\_unfitw/} are run from the
same directory as
\texttt{PYTHONPATH=. venv/bin/python scribe\_layout\_v2/runs/compare\_unfitw/<script>.py [args]}:
\begin{center}\small
\begin{tabular}{p{5.3cm}p{9.9cm}}
\toprule
script (arguments) & output and use\\
\midrule
\texttt{compare\_5seed.py} & plan-level C2ST, FFT C2ST, residual-scale $W_1$, test, 5 seeds $\to$ \texttt{compare\_5seed.txt} (Table~\ref{tab:plan})\\
\texttt{compare\_5seed\_dc.py} & the same for the deconvolved model $\to$ \texttt{compare\_5seed\_dc.txt} (Table~\ref{tab:plan})\\
\texttt{compare\_5seed\_val.py} & the same on the validation split $\to$ \texttt{compare\_5seed\_val.txt} (Table~\ref{tab:val})\\
\texttt{compare\_5seed\_wh\_test.py}, \texttt{compare\_5seed\_wh\_val.py} & default vs Huber-aware per-word correction, test and val, 5 seeds $\to$ \texttt{compare\_5seed\_wh\_test.txt}, \texttt{compare\_5seed\_wh\_val.txt} (Table~\ref{tab:huber})\\
\texttt{fft\_wobble\_scale.py} & FFT C2ST vs drawn-wobble scale, 3 seeds $\to$ \texttt{fft\_wobble\_scale.txt} (\S\ref{sec:fftdiag})\\
\texttt{leverage\_sim.py} & per-word wobble recovery $\to$ \texttt{leverage\_sim.txt} (Table~\ref{tab:levsim})\\
\texttt{curvnoise2.py} & curvature noise decomposition by $n$, train (Table~\ref{tab:curvnoise}); prints only\\
\texttt{sag.py} (\textit{model.pt}) & line sag, curvature sd and adjacent steps, test seed 0 (Table~\ref{tab:sag}); prints only\\
\texttt{drift.py} & plan-level slope, drift, curvature and pitch $W_1$, 5 seeds (Table~\ref{tab:drift}); prints only\\
\texttt{split\_diag.py} (\textit{run\_dir method}) & extracted vs drawn lines: split, single-piece, extra fragments (\S\ref{sec:e2e}); prints only\\
\texttt{split\_why.py} (\textit{run\_dir method}) & anchor steps and merge residuals at split points; (\textit{mid}$-$headline) spread (\S\ref{sec:dc}); prints only\\
\texttt{surv.py} (\textit{run\_dir}) & drawn slope and bow of split vs kept v2 lines (\S\ref{sec:e2e}); prints only\\
\bottomrule
\end{tabular}
\end{center}

\begin{center}\small
\begin{tabular}{p{7.2cm}p{8.0cm}}
\toprule
\texttt{extract.py}, \texttt{anchor\_noise.py} & combined anchor; noise deconvolution\\
\texttt{unfit.py} & line-scale table, per-word leverage correction, Huber-aware factor (not used)\\
\texttt{features.curv\_noise\_var} & curvature noise variance $v_i$\\
\texttt{nn\_utils.MDN.nll}, \texttt{model.head\_nll} & convolved MDN likelihood (\texttt{deconv})\\
\texttt{line\_clearance.py} & real clearance thresholds\\
\texttt{fit\_tilt.py} & rotation gain $\beta$\\
\texttt{features.py}, \texttt{model.py}, \texttt{train.py}, \texttt{api.py} & MDN chain, history and latent styles\\
\texttt{bscribe\_layout.py}, \texttt{relayout.py}, \texttt{bscribe\_run.py} & BengaliScribe integration (planner reset, margin guard)\\
\texttt{tests/} & 70 unit tests\\
\texttt{evaluate.py}, \texttt{baselines.py}, \texttt{e2e\_check.py} & plan-level and end-to-end evaluation\\
\texttt{scribe\_layout\_fft2d/}, \texttt{scribe\_layout\_fourier/} & 2-D FFT and 1-D Fourier experiments\\
\bottomrule
\end{tabular}
\end{center}

\begin{thebibliography}{9}
\bibitem{bovy2011}
J.~Bovy, D.~W. Hogg and S.~T. Roweis (2011).
Extreme deconvolution: inferring complete distribution functions from noisy,
heterogeneous and incomplete observations.
\emph{Annals of Applied Statistics} 5(2B), 1657--1677.
\end{thebibliography}

\end{document}
